\chapter{Arithmetic fusion as singular interfaces}
\label{v4-arith:asi:chapter}
\chaptermark{Arithmetic fusion as singular interfaces}

The marked network of Chapter~\ref{v4-arith:mf:chapter} has value
\(b_g\), and its specified duality closes it to \(-e_g\).
Insert the same junctions into two Heisenberg currents. Their
contraction produces the double pole
\[
 -\frac{e_g}{(z_1-z_3)^2}.
\]
The intervening insertions remain part of the ordered coefficient
product. Thus the singular amplitude retains a marked fusion
composition that its scalar residue alone would not determine.

We construct this realization on all oscillator states and every
collision stratum. It has infinite-dimensional coefficient spaces.
A second realization uses the finite bimodules of the preceding
chapter, together with the categorical symmetry that singles out
their allowed maps. Tensoring them with the \(E_8\) lattice theory
gives finite conformal grades and convergent relative sewing.

For the order-two arithmetic category, two further realizations have
a single conformal vacuum in their ambient vertex algebra.
Changing the quadratic embedding gives the \(E_6\) module category
at affine \(\mathfrak{sl}_2\) level ten. The original embedding is
realized by the Virasoro theory of central charge \(-39/10\).
We compare their actual junctions, full descendant fields and
collision coefficients. The coefficient algebra, the finite
bimodule, the conformal state space and the conformal-block bundle
will remain separate objects throughout.

\section{The junction and its coefficient field}

Use the arithmetic data of Section~\ref{v4-arith:mf:arithmetic} and write
\(F=F_I\), \(B=\mathcal B_I\), \(\mathcal C=\mathcal C_I\) and
\(G=G_I\). Thus \(F\) has characteristic zero and
\[
 G=(\varepsilon-1)^{-1}I/I,\qquad
 N=|G|=\varepsilon+\varepsilon^{-1}-2.
\]

The generators \(X_g,X'_g,T_g,T'_g\), their endomorphisms
\(\alpha_g,\rho\), and the elements \(U_g\) are those of
Section~\ref{v4-arith:mf:leavitt}. Primes denote algebraic generators,
not adjoints. In the notation used below,
\begin{equation}\label{v4-arith:asi:normalizations}
 e_g=E_g/E_0,\quad b_g=E_g/\sqrt N,\quad
 b_0=\frac{\varepsilon-1}{N},\quad
 d=-b_0^{-1}=\varepsilon^{-1}-1.
\end{equation}

The positive roots obey
\(\sqrt N=\sqrt\varepsilon-\varepsilon^{-1/2}\).
Put \(\beta=\sum_gb_g\), \(a=\beta/(Nb_0)\) and
\(c=b_0/\beta\). Here \(c\) is the coefficient denoted \(c_1\) in \eqref{v4-arith:mf:parameters}. The Gaussian \(Q\) and bicharacter \(\chi\)
retain the conventions of \eqref{v4-arith:mf:chi}.

Tensor categories and their functors are linear over their stated
fields, except where semilinearity is specified. Complexes are
cohomologically \(\mathbb Z\)-graded, with the Koszul sign for tensor
products of homogeneous maps.

The endomorphism presentation gives
\begin{equation}\label{v4-arith:asi:fusion-identities}
 \alpha_g\rho=\rho,\qquad
 \rho\alpha_g=\operatorname{Ad}(U_g)\rho,\qquad
 \alpha_g(U_h)=\chi(g,h)U_h.
\end{equation}
Here \(\operatorname{Ad}(u)(z)=uzu^{-1}\). The first two identities
are proved in Theorem~\ref{v4-arith:mf:fusion}. The last follows by
applying \(\alpha_g\) to \eqref{v4-arith:mf:U}.
The specified cup and cap are \(X_0\) and \(dX'_0\).

The fusion ring records the quadratic period. For
\(r=[\rho]\) and \(S=\sum_g[\alpha_g]\), the relations
\(r^2=Nr+S\) and \(rS=Nr\) give
\[
 (1+r)(1-r+S)=1.
\]
The positive dimension character takes \(1+r\) to
\(\varepsilon\), while the specified duality dimension takes it to
\(\varepsilon^{-1}\). The inverse is a virtual class, not an inverse
object for \(1\oplus\rho\).

The marked network of Proposition~\ref{v4-arith:mf:network-value} is
\begin{equation}\label{v4-arith:asi:network}
 \mathcal N_g=T'_0T'_0\rho(T_g)T_0.
\end{equation}
Its value and its duality closure are
\begin{equation}\label{v4-arith:asi:closed-network}
 \mathcal N_g=b_g1,\qquad
 dX'_0\rho(\mathcal N_g)X_0=db_g=-e_g.
\end{equation}
The calculation following \eqref{v4-arith:mf:closure} shows why this closure must
retain the marked object: in the ambient algebra,
\begin{equation}\label{v4-arith:asi:HH-unit}
 \sum_g\bigl([X'_g,X_g]+[T'_g,T_g]\bigr)=(2N-1)1,
\end{equation}
so the unit has zero Hochschild class.

For \(\varepsilon=2+\sqrt3\), the simples are \(1,g,\rho\), with
\(g^2=1\), \(g\rho=\rho g=\rho\) and
\(\rho^2=1\oplus g\oplus\rho^{\oplus2}\). Their coefficients are
\[
 b_0=\frac{1+\sqrt3}{2},\qquad b_1=\frac{1+i}{2},\qquad d=1-\sqrt3.
\]
The second marked network closes to the negative of
\begin{equation}\label{v4-arith:asi:example-unit}
 e_1=\frac{(\sqrt3-1)(1+i)}2,
 \qquad e_1^4+2e_1^3+2e_1^2-2e_1+1=0.
\end{equation}
The singular construction will recover this value as the coefficient
of a double pole.

\section{Operator coefficients at a nonconsecutive collision}

A coefficient multiplying a field may be noncommutative. This changes the
operation at a collision that passes an intervening insertion. For example,
take $A=C=E_{12}$ and $D=E_{21}$ in $\operatorname{Mat}_2(F)$. Then
$AC=0$ while $ADC=E_{12}$. The algebra $B$ contains such matrix units:
choose two of its positive generators $L_1,L_2$, and put
$E_{ij}=L_iL'_j$. This gives a matrix corner with identity $E_{11}+E_{22}$.

The exchange of Section~\ref{v4-sec:regular-contextual-exchange}
retains this context. It is the linear map
\begin{equation}\label{v4-arith:asi:contextual-crossing}
 c_B:B\otimes B\longrightarrow B\otimes B,\qquad
 c_B(x,y)=xy\otimes1+1\otimes xy-x\otimes y.
\end{equation}
Let $m$ denote multiplication in $B$. Direct expansion gives
\begin{gather}
 c_B^2=1,\quad mc_B=m,\quad
 c_B(1,x)=(x,1),\quad c_B(x,1)=(1,x),\notag\\
 (c_B)_{12}(c_B)_{23}(c_B)_{12}
 =(c_B)_{23}(c_B)_{12}(c_B)_{23},\label{v4-arith:asi:coefficient-relations}\\
 c_B(1\otimes m)=(m\otimes1)(1\otimes c_B)(c_B\otimes1),\notag\\
 c_B(m\otimes1)=(1\otimes m)(c_B\otimes1)(1\otimes c_B).\notag
\end{gather}
For example, the two braid words on $x\otimes y\otimes z$ both give
\begin{align*}
 xyz\otimes1\otimes1-xy\otimes z\otimes1+1\otimes1\otimes xyz
 +1\otimes xy\otimes z\\
 {}+x\otimes yz\otimes1-1\otimes x\otimes yz-x\otimes y\otimes z.
\end{align*}
No dimension hypothesis occurs in these identities.

Let $f$ be a map between finite ordered sets. Define $q_f$ by three steps.
Sort the inputs into the fibers of $f$ with adjacent $c_B$'s.
Multiply inside each fiber, and insert $1_B$ in an empty fiber. The relations
\eqref{v4-arith:asi:coefficient-relations} imply
\begin{equation}\label{v4-arith:asi:q-functor}
 q_gq_f=q_{gf},\qquad q_{\mathrm{id}}=1.
\end{equation}

To prove this, use the braid and involution relations to change a sequence
of adjacent crossings. Crossings inside one fiber disappear by $mc_B=m$. The last two
relations move a merger through a crossing. Every composite therefore
reduces to the sorting diagram for the composite map, followed by the
multiplications in its fibers. Associativity and the unit equations remove
the remaining choices. These moves also prove naturality in every unital
algebra homomorphism $B\to B'$.

At three inputs, with the remaining blocks ordered as $2|13$, this gives
\begin{equation}\label{v4-arith:asi:q13}
 q_{2|13}(x,y,z)=xy\otimes z+1\otimes xyz-x\otimes yz,
 \qquad m q_{2|13}(x,y,z)=xyz.
\end{equation}
At four inputs,
\[
 q_{13|24}(x,y,z,w)=xyz\otimes w+x\otimes yzw-xy\otimes zw.
\]
The intermediate factors have not been replaced by the product of the
colliding entries. This is the coefficient information needed for the
geometric residue identities.

\section{A singular construction in every arity}

Use the level-one Heisenberg vertex algebra \(V\)
of Section~\ref{v4-sec:singular-heisenberg}, now over \(F\). Its state space is
$F[x_1,x_2,\ldots]$, its vacuum is $1$, and its current modes are
\[
 J_{-m}=x_m,\qquad J_m=m\partial_{x_m}\quad(m>0),\qquad J_0=0.
\]

Write $J=x_1$ for the current state. The translation operator is
$T_V=\sum_{m>0}m x_{m+1}\partial_{x_m}$.
The field of $x_m$ is $\partial_z^{m-1}J(z)/(m-1)!$. The field of a
monomial is the normal ordered product of these fields. The commutator
$[J_m,J_n]=m\delta_{m+n,0}$ gives
\begin{equation}\label{v4-arith:asi:Heisenberg-pole}
 J(z)J(w)\sim\frac1{(z-w)^2}.
\end{equation}

The double pole is the contraction
$\sum_{m>0}m z^{-m-1}w^{m-1}=(z-w)^{-2}$ in the region $|z|>|w|$.
Thus the pole and its normalization are fixed by the oscillator
commutator. The arithmetic coefficients will multiply this contraction.
All states are polynomials, and all products have an input-dependent finite
pole order. No completion of the coefficient algebra is imposed.

Take $X=\mathbb A^1_F$ and $H=\omega_X\otimes_F V$. We use right
differential modules, with
\[
 (f(z)\,dz\otimes v)\partial_z
 =-f'(z)\,dz\otimes v-f(z)\,dz\otimes T_Vv.
\]
For $t=z_1-z_2$, the local-current face is the full principal-part map
\[
 r_V(f;v,w)=\operatorname{pp}_{t=0}\bigl(fY_V(v,t)w\bigr).
\]
Normal derivatives are retained: $t^{-q-1}$ corresponds to
$\delta\partial_t^q/q!$, where $\delta t=0$. Thus multiplication by $t$
lowers the principal-part order and right multiplication by $\partial_t$
raises it with coefficient $q+1$.

We use the local-current Chevalley--Cousin construction of
Beilinson--Drinfeld, \S3.4.11. Its input is the chiral differential module of $V$,
with the above full faces. The two-step face sum vanishes: disjoint
collisions cancel with their determinant signs, and a three-block
collision is the chiral Jacobi identity. For the Heisenberg algebra these
identities also follow from Wick expansion and the central commutator of
the modes. This local-current input is distinct from the ordered
coefficient multiplication in $B$.

For $n$ labeled positions, the ordered collision scheme $R_n$ has one
chart $X^n$ for each ordering. Identify two charts away from the diagonals
on which their orderings disagree.

A partition $\pi$ has
one coordinate for each block. Write $i_\pi:X^\pi\to X^n$ for its diagonal
and $j_\pi$ for the distinct-block open subset of $X^\pi$. Define
\[
 Q_\pi(V)=(i_\pi)_*(j_\pi)_*j_\pi^*
       \left(\boxtimes_{A\in\pi}H\right).
\]

Order the blocks by their last entries and retain the later root at a
merger. Orient its normal coordinate by $z_i-z_j$ for $i<j$.
With the ordered density line the face that removes block $i$ has sign
$(-1)^{i-1}$. These conventions fix the signed current face
$d^V_{\pi,\eta}$ for every elementary coarsening $\pi\lessdot\eta$.

The associative construction uses consecutive cuts in a fixed word,
meaning a finite ordered list of insertion positions.
To specify its carrier, let $\operatorname{OrdSurj}_+$ have the nonempty
finite ordered sets as objects and monotone surjections as arrows. Put
\[
 \mathcal W_X=D(F)\times
 \varprojlim_{\operatorname{OrdSurj}_+}D\operatorname{Mod}(X^n),
\]

The limit is taken in stable $\infty$-categories. We use unbounded right
differential modules with quasi-coherent underlying modules.
For $\alpha:[n]\twoheadrightarrow[m]$, the transition is
extraordinary inverse image along the diagonal
$d_\alpha(x)_i=x_{\alpha(i)}$. An object includes the transition
identifications and their coherences. The factor $D(F)$ records the empty
configuration independently.

For $0<p<n$, let $U_p\subset X^n$ be the open set where the first $p$
positions are distinct from the remaining positions.
Write $j_p:U_p\hookrightarrow X^n$. The cut product is
\[
 (M\star N)_n=
 \bigoplus_{p=0}^n(j_p)_*
       \left.(M_p\boxtimes N_{n-p})\right|_{U_p}.
\]

An empty block means scalar tensoring on all of $X^n$. The unit is $F$
at the empty configuration and zero elsewhere. A strong cut coalgebra
has a counit and coassociative cut maps whose restrictions to each $U_p$
are isomorphisms. For a connected coaugmented such coalgebra $C$, its cut
cobar is the tensor algebra $T_\star(C_+[-1])$ on its positive part.
Writing $tc$ for desuspension, its differential is the derivation with
\[
 d(tc)=-t(d_Cc)+
 \sum_{\overline\Delta c=c'\otimes c''}
             (-1)^{|c'|+1}(tc')(tc'').
\]
Here $\overline\Delta$ is the sum of the nonempty cut maps.

\begin{theorem}\label{v4-arith:asi:all-arity}
Let $B$ be a nonzero unital associative $F$-algebra concentrated in degree
zero. The chartwise formula
\begin{equation}\label{v4-arith:asi:coefficient-complex}
 C_n(V,B)=\bigoplus_{\pi\vdash[n]}Q_\pi(V)[|\pi|]\otimes_F B^{\otimes|\pi|},
 \qquad d=\sum_{\pi\lessdot\eta}d^V_{\pi,\eta}\otimes q_{\pi,\eta}
\end{equation}
defines a complex on every ordered chart and glues on $R_n$.
The complexes on the standard word charts form a connected strong cut
coalgebra in $\mathcal W_X$, with $C_0(V,B)=F$. This coalgebra is
natural in unital algebra maps. Its cut cobar
\[
 \mathcal A_B=\Omega_{\mathrm{cut}}C(V,B)
\]
is an associative algebra in the ordered cut-convolution category.
Its binary field system is
\begin{equation}\label{v4-arith:asi:singular-fields}
 Y(v\otimes b,z)(w\otimes c)=Y_V(v,z)w\otimes bc,
 \qquad T=T_V\otimes1.
\end{equation}
\end{theorem}

\begin{proof}
For fixed $\pi$ and $\nu$ separated by two mergers,
\eqref{v4-arith:asi:q-functor} makes every coefficient composite equal to
$q_{\pi,\nu}$. The corresponding component of $d^2$ is therefore
\[
 \left(\sum_{\pi\lessdot\eta\lessdot\nu}
 d^V_{\eta,\nu}d^V_{\pi,\eta}\right)\otimes q_{\pi,\nu}=0.
\]
On a chart overlap use the usual current permutation and determinant
identification, and use $c_B$ to reorder coefficient blocks. The braid
relation gives the overlap cocycle, and \eqref{v4-arith:asi:q-functor} gives
compatibility with the differential. A partition containing an internal
pair whose orders disagree has empty support on that overlap.

Under an ordered collision $\alpha:[n]\twoheadrightarrow[m]$, a partition
term survives exceptional pullback only when it is saturated for $\alpha$.
Otherwise its distinct-block open factor is pulled to a collision, and
$i^!j_*=0$ makes that term vanish. The surviving terms are exactly
$\pi=\alpha^{-1}(\eta)$, with the same coefficient factors and faces as
$C_m(V,B)$. Closed-diagonal base change identifies their differential
modules. For composite collisions the same saturated partitions survive,
so the comparisons compose.

On the open subset where the first $p$ positions are distinct from the
last $n-p$, every partition crossing that cut disappears. A remaining
partition is uniquely a partition of each block. Its coefficient complex
is the external product of the two block complexes. These are the strong
cut comparisons, and their associativity follows from splitting into three
blocks.

The cobar is the tensor algebra on the shifted positive
coefficient, with differential given by $d$ and these cut maps. At a fixed
configuration its positive tensor length is bounded by the arity. No
infinite sum in that length is needed.

All arguments take place in the category of differential modules allowing
arbitrary coefficient vector spaces. Tensoring a differential module by
a vector space means the corresponding direct sum. Localization and
closed-diagonal direct image commute with these sums. Thus neither an
algebraic dual of $B$ nor a finite-rank Riemann--Hilbert assertion is used.

For a finitely generated algebra, filter by word length
and put the total word-length filtration on $B^{\otimes s}$. The maps
$c_B$ and $q_f$ preserve it. Each finite-arity coefficient diagram is the
filtered colimit of its finite-dimensional coefficient diagrams. This
applies to the Leavitt algebra, even though its finite word-length spaces
are not subalgebras. Naturality and the binary assertion follow directly
from the formulas.
\end{proof}

The positive-cut contraction of
Section~\ref{v4-sec:bar-positive} retains every nonconsecutive stratum.
Here its normal derivatives and coefficient maps can be written
explicitly. The cobar retracts onto a complex supported on
\[
 Z_n=\bigcap_{p=1}^{n-1}\ \bigcup_{i\le p<j}\Delta_{ij};
 \qquad Z_3=\Delta_{13},\quad
 Z_4=\Delta_{14}\cup(\Delta_{13}\cap\Delta_{24}).
\]

The differential modules on these strata retain all normal derivatives.
For a partition $\pi$, let
$J(\pi)$ be the cuts crossed by none of its blocks. A cobar term is indexed
by a subset $Q\subseteq J(\pi)$, with coefficient
\[
 \bigl(Q_\pi(V)\otimes B^{\otimes|\pi|}\bigr)
             [|\pi|-|Q|-1].
\]

If $r_{\le c}$ is the number of blocks to the left of a cut $c$, adding
$c\notin Q$ has sign
$(-1)^{\#\{q\in Q:q<c\}+1-r_{\le c}}$.
Merging root-indexed blocks $i<j$ in the $\ell$th consecutive segment
has sign $(-1)^{\ell-1+i}$ and coefficient map
$r^V_{ij}\otimes q_{ij}$, where $r^V_{ij}$ is the unsigned current face.
Write $v$ and $\delta$ for these two parts of the differential.

For $J(\pi)\ne\varnothing$, put $c_0=\min J(\pi)$. Define $h_0$ to remove
$c_0$ with sign $(-1)^{1-r_{\le c_0}}$ when $c_0\in Q$, and to be zero
otherwise. It is also zero when $J(\pi)=\varnothing$.

Let $i$ include the terms with no admissible cut, and let $p_0$ project
onto them. Adding and removing $c_0$ gives $vh_0+h_0v=1-ip_0$.
All other cut terms cancel in pairs.
Merging blocks cannot create an admissible cut, so these terms form a
subcomplex for $\delta$. The perturbed projection and homotopy are
\[
 p=p_0\sum_{q=0}^{n-1}(-\delta h_0)^q,\qquad
 H=h_0\sum_{q=0}^{n-1}(-\delta h_0)^q.
\]

Each residue lowers the number of blocks, which proves finiteness. With
these maps one has $dH+Hd=1-ip$.
After suspension, put $\bar H=-sHs^{-1}$, $\bar i=sis^{-1}$ and
$\bar p=sps^{-1}$. If $b_2$ is the suspended cobar product, its transferred
operations are specified in every arity by
\[
 F_1=\bar i,\quad
 \Lambda_n=\sum_{a+b=n}b_2(F_a\otimes F_b),\quad
 F_n=-\bar H\Lambda_n,\quad B_n=\bar p\Lambda_n\quad(n\ge2).
\]
These are finite homological perturbation formulas. The coefficient maps
commute with the contraction, so algebra homomorphisms preserve the full
transferred system.

For the arithmetic algebra, \eqref{v4-arith:asi:singular-fields} realizes
\eqref{v4-arith:asi:network} with a genuine pole. Let $\epsilon_V$ extract the
constant Fock coefficient. It is a linear functional, not a character of
the vertex algebra. In the expansion region $|z_1|>|z_3|$,
\begin{align}
 &(\epsilon_V\otimes1)
 Y(J\otimes T'_0,z_1)Y(1\otimes T'_0,z_2)
 Y(J\otimes\rho(T_g),z_3)Y(1\otimes T_0,z_4)(1\otimes1)
 \notag\\[-1mm]
 &\hspace{30mm}=\frac{b_g}{(z_1-z_3)^2}\,1_B.
 \label{v4-arith:asi:singular-network}
\end{align}

Closing the coefficient line multiplies this by $d$, giving
$-e_g/(z_1-z_3)^2$. The rational function determines the other expansion
regions. Arithmetic enters through the marked endomorphisms and junctions.
The Heisenberg level remains one.

\section{Relative interfaces and the connected comparison}

The field realization retains compositions of coefficients. For its
tensor structure, recall the bimodules of
Section~\ref{v4-sec:endomorphism-interfaces} with the following conventions.
Put $R=B^{\mathrm{op}}$. For an endomorphism $\sigma$ of $B$, let
$P_\sigma$ have underlying vector space $B$ and actions
\[
 a^{\mathrm{op}}\cdot b\cdot c^{\mathrm{op}}=\sigma(c)ba.
\]

An intertwiner $x:\sigma\to\tau$ acts by $b\mapsto xb$.
Every left $R$-linear map has this form, and right $R$-linearity is exactly
the intertwining equation. Hence this construction is fully faithful.
The tensor comparison is
\begin{equation}\label{v4-arith:asi:bimodule-product}
 P_\sigma\otimes_R P_\tau\longrightarrow P_{\sigma\tau},\qquad
 b\otimes c\longmapsto\sigma(c)b.
\end{equation}

It is balanced, has inverse $z\mapsto z\otimes1$, and is associative
because $\sigma(\tau(e)c)b=\sigma\tau(e)\sigma(c)b$.
It takes tensor products of intertwiners to $x\sigma(y)$.
The unit is $P_{\mathrm{id}}=R$. Every $P_\tau$ is free as a left
$R$-module, so the same formula computes the derived balanced product.

Let $A=\mathcal A_F$ be the Heisenberg cut cobar, let $\star$ denote
ordered cut convolution, and put
\[
 \mathfrak R=A\otimes_F R,\qquad
 \mathfrak P_\sigma=A\otimes_F P_\sigma.
\]
Coefficient vector spaces act centrally on differential modules. The
$A$-factors are always multiplied in their original order.
\begin{theorem}\label{v4-arith:asi:relative-realization}
There is a faithful strong tensor functor
\[
 \mathcal C\longrightarrow
 \operatorname{Bimod}_{\mathfrak R}(\mathcal W_X),\qquad
 \sigma\longmapsto\mathfrak P_\sigma,
\]
where $\mathcal W_X$ is the derived ordered cut-convolution category and
the target tensor product is $\star^{\mathbb L}_{\mathfrak R}$.
It preserves the specified duality and all fusion-network evaluations.
\end{theorem}
\begin{proof}
For coefficient bimodules of this form, the two-sided bar gives
\begin{equation}\label{v4-arith:asi:relative-bar}
 (A\otimes P)\star^{\mathbb L}_{A\otimes R}(A\otimes Q)
 \simeq A\otimes(P\otimes_R^{\mathbb L}Q).
\end{equation}
The comparison multiplies the $A$-factors in simplicial degree $n$ in
their original order and leaves the coefficient bar
$P\otimes R^{\otimes n}\otimes Q$.
One can compute using the free left coefficient resolution of $Q$.
After extension by $A$ its terms are free left $A\otimes R$-modules.
The augmented resolution contracts by its extra unit degeneracy.
This proves \eqref{v4-arith:asi:relative-bar}.

Equivalently, for $P_\sigma$
and $P_\tau$ it is the internal version of the explicit free-module
isomorphism \eqref{v4-arith:asi:bimodule-product}.

All parenthesizations for three or more factors realize the same
multisimplicial bar. Multiplying the geometric factors in order gives
the same map in each case, proving the associativity and unit coherences.

Combine this comparison with \eqref{v4-arith:asi:bimodule-product} and extend
to finite sums and retracts. Evaluation at the empty configuration
recovers $P_\sigma$ and its coefficient maps, proving faithfulness.
The images of the cup and cap obey the same snake equations. Applying
the functor to \eqref{v4-arith:asi:closed-network} proves the assertion about
network values.
\end{proof}

Here \(\mathfrak R_\emptyset=R\). This empty algebra is part of
the tensor comparison. There is also a connected construction
with empty coefficient $F$, and the relation between them is an actual
algebra map. Apply Theorem~\ref{v4-arith:asi:all-arity} to $R$, and set
$G_R=\mathcal A_R$. If $m_\pi:R^{\otimes|\pi|}\to R$ multiplies in
block order, then
\[
 m_\eta q_{\pi,\eta}=m_\pi.
\]

Let $\tau:C(V,F)_+\to A$ be the universal cobar twisting map. On the
$\pi$-term define $t_R=\tau\otimes m_\pi$. The last identity makes it
compatible with residues, and on a cut the total product is the product
of the two block products. Thus $t_R$ obeys the twisting equation and
induces
\begin{equation}\label{v4-arith:asi:theta}
 \theta_R:G_R\longrightarrow\mathfrak R.
\end{equation}

On a cobar monomial it retains the geometric monomial and multiplies
all coefficient factors in their original order. This construction takes
place in the word category $\mathcal W_X$ and is natural in $R$.
Its empty component is the unit map $F\to R$.

Intertwiners have a linear connected realization as well. For an
$R$-bimodule $P$, form the square-zero algebra $R\ltimes P$.
Give $R$ weight zero and $P$ weight one. Multiplication and the exchange
preserve total weight, so $\mathcal A_{R\ltimes P}$ is weight graded.
Its weight zero part is $G_R$. Its weight one part, denoted
$L_R(P)$, is a $G_R$-bimodule.

Each coefficient in weight one contains
exactly one $P$-factor. Consequently $P\mapsto L_R(P)$ is additive and
linear on bimodule maps. Its one-point component is $H\otimes P$, so
it is faithful on degree-zero coefficient maps.

Taking weight one in
\eqref{v4-arith:asi:theta} for $R\ltimes P$ gives the comparison
\begin{equation}\label{v4-arith:asi:tangent-comparison}
 L_R(P)\longrightarrow\operatorname{Res}_{\theta_R}(A\otimes P).
\end{equation}
It is the identity at one point. At the empty configuration it is $0\to P$.

This extraction of weight one is necessary. For example, with
$M_\sigma={}_\sigma B_B$, an intertwiner $x$ gives an algebra map
\[
 B\ltimes M_\sigma\longrightarrow B\ltimes M_\tau,\qquad
 (b,m)\longmapsto(b,xm).
\]
The algebra map has the identity on its weight zero part, so it is not
linear as a function of $x$. Its weight one map is linear. Such maps
realize the intertwiners, but do not identify their transported tensor
product with $\star^{\mathbb L}_{G_R}$.

\section{Analytic states and relative sewing}
\label{v4-arith:asi:sec:analytic}

Fix an embedding $F\hookrightarrow\mathbb C$ and extend scalars along it.
The oscillator construction has an analytic realization on all states.
This also gives an analytic meaning to the relative interfaces.

For an $R$-bimodule $P$ put
\[
 \mathsf H(P)=V\otimes P,\qquad
 \mathsf H(P)_n=V_n\otimes P,\qquad
 \widehat{\mathsf H(P)}=\prod_{n\geq0}\mathsf H(P)_n.
\]
Here $\deg x_m=m$. States belong to the direct sum. Outputs of fields
belong to the product. The restricted dual is
$\bigoplus_{n\geq0}\mathsf H(P)_n^*$, with the full algebraic dual in each
degree.

The conformal vector is $\frac12x_1^2\otimes1$ in
$\mathsf A=\mathsf H(R)$. Its central charge is one.

An \emph{oscillator interface} is $\mathsf H(P)$ with its Heisenberg
action and commuting left and right $R$-actions. Morphisms commute with
these three actions. If $f:P\otimes_RQ\to S$ is a bimodule map, define
\begin{equation}\label{v4-arith:asi:interface-field}
 \mathcal I_f(v\otimes p,z)(w\otimes q)
 =Y_V(v,z)w\otimes f(p\otimes_Rq).
\end{equation}
A balanced intertwining operator means an intertwining operator for the
Heisenberg vertex algebra, with these left, right, and balancing identities
on its coefficient factors. This definition includes every oscillator
descendant.

\begin{theorem}\label{v4-arith:asi:analytic-relative}
The oscillator interfaces have the following properties.
\begin{enumerate}[label=\textup{(\roman*)}]
\item Their morphisms are exactly $1_V\otimes f$, for bimodule maps $f$.
Every balanced intertwining operator is uniquely of the form
\eqref{v4-arith:asi:interface-field}. Consequently their relative fusion is
represented by
\[
 \mathsf H(P)\boxtimes_{\mathsf A}\mathsf H(Q)
 =\mathsf H(P\otimes_RQ).
\]
\item Products and iterates of their fields converge absolutely after
pairing with any restricted-dual state, in the usual nested-disc regions.
They are expansions of rational functions. Every genus-zero sewing of
these interfaces respects their tensor associativity and units.
\item The functor $\sigma\mapsto\mathsf H(P_\sigma)$ is fully faithful
and strong tensor. It preserves the specified cups, caps, and marked
arithmetic junctions on all states. Its tensor-product field is
\begin{equation}\label{v4-arith:asi:full-tensor-field}
 \mathcal I_{\sigma,\tau}(v\otimes b,z)(w\otimes c)
 =Y_V(v,z)w\otimes\sigma(c)b.
\end{equation}
\item Taking full normal principal parts gives the relative chiral
interfaces of Theorem~\ref{v4-arith:asi:relative-realization}. The comparison
commutes with every collision, consecutive cut, and cobar differential.
\end{enumerate}
\end{theorem}

\begin{proof}
The vectors killed by all positive Heisenberg modes in $\mathsf H(P)$
are exactly $1\otimes P$. A map commuting with the modes is determined
there, because the negative modes generate every polynomial state.
The coefficient actions give the bimodule condition. This proves the
morphism assertion.

For an intertwining operator $\mathcal I$, translation gives
$\partial_z\mathcal I(1\otimes p,z)=0$. Its value on $1\otimes q$ has
weight zero, hence is $1\otimes f(p,q)$. The Heisenberg commutator and
iterate identities determine its values on all states by induction on
the number of negative modes. They give exactly
\eqref{v4-arith:asi:interface-field}. Balancing says
$f(pr,q)=f(p,rq)$, and the outer actions say that the induced map on
$P\otimes_RQ$ is a bimodule map. This proves the universal property in
\textup{(i)}.

Wick expansion of finitely many polynomial states has finitely many
contraction patterns. Each contraction is a derivative of
$(z_i-z_j)^{-2}$. Pairing with a restricted-dual state leaves a finite
sum of rational functions. Their Taylor expansions converge absolutely
when each inner disc is contained in its outer disc and avoids the
remaining insertions. The finitely many coefficient products lie in a
finite-dimensional subspace of the algebraic coefficient space, so no
norm or completion of $R$ is needed. Changing a tree changes only the
expansion region of the same function and the bracketing of the same
balanced coefficient product.

This proves \textup{(ii)}.

Apply \textup{(i)} to \eqref{v4-arith:asi:bimodule-product}. Its associativity,
unit, and intertwiner equations prove \textup{(iii)}, first on
endomorphisms and then on finite sums and retracts. The cups and caps
are the corresponding coefficient maps tensored with $1_V$. Their
snake identities therefore hold on every descendant.

Finally, the state identification sends
$x_{m_1}\cdots x_{m_l}\otimes p$ to
$J_{-m_1}\cdots J_{-m_l}1\otimes p$.
It intertwines translation and all fields. It therefore identifies the
right differential-module actions and every full principal part.
The coarsenings multiply the same ordered coefficients on both sides.
The cut maps select the same partitions. These identities identify the
whole Cousin diagram and hence its cobar, proving \textup{(iv)}.
\end{proof}

The sum over oscillator descendants in a sewing is explicit. For a
partition $\lambda=(1^{m_1}2^{m_2}\cdots)$ put
$x_\lambda=\prod_jx_j^{m_j}$ and
$z_\lambda=\prod_jj^{m_j}m_j!$. Normalize the Heisenberg pairing by
$\langle1,1\rangle=1$ and
$\langle J_{-m}v,w\rangle=\langle v,J_mw\rangle$.
Identify the opposite-coordinate dual Fock space through this pairing.

Its sewing tensor is
\begin{equation}\label{v4-arith:asi:oscillator-sewing}
 \mathcal P(t)=\sum_\lambda
       \frac{t^{|\lambda|}}{z_\lambda}\,x_\lambda\otimes x_\lambda,
 \qquad |\lambda|=\sum_jjm_j.
\end{equation}
Coefficient factors are contracted over $R$, by the relative evaluation.
They are not summed over a vector-space basis of $R$. The Wick argument
in the proof establishes the convergence of the resulting matrix
coefficients on a sewing annulus. Omitting the positive oscillator
degrees would change, for example, the expansion
\[
 \frac1{(z-w)^2}=\sum_{n\geq1}n\,z^{-n-1}w^{n-1}.
\]

Thus every arithmetic category in the construction has a full analytic
ordered realization, with nonconstant fields and all descendants.
It is a relative theory with infinite-dimensional weight-zero
coefficients. This qualification cannot be removed by truncating the
same coefficient action. A finite-dimensional unital $B$-module $W$
would satisfy $W\simeq W^{\oplus2N}$ by the Leavitt generators, hence
$\dim W=2N\dim W$ and $W=0$. A realization with finite-dimensional
homogeneous state spaces must use a different coefficient carrier.

\section{Finite coefficients and categorical symmetry}
\label{v4-arith:asi:sec:finite-coefficients}

Semisimplicity provides a finite carrier. Its extra structure is essential:
the allowed maps must respect the right action of the fusion category.
Without this condition a finite matrix realization introduces new junctions.

Continue over $\mathbb C$. Choose representatives $(S_i)_{i\in\mathcal I}$ of
the simple objects of $\mathcal C$, with $S_0=1$, and put
\[
 R_0=\bigoplus_{i\in\mathcal I}\mathbb C e_i,\qquad
 e_i\mathsf P_Xe_j=\operatorname{Hom}_{\mathcal C}(S_i,X\otimes S_j).
\]
The $e_i$ are orthogonal idempotents with sum one. A map $x:X\to Y$
acts on $\mathsf P_X$ by $f\mapsto(x\otimes1)f$. With
$a_{X,Y,Z}:(X\otimes Y)\otimes Z\to X\otimes(Y\otimes Z)$,
the tensor comparison is
\begin{equation}\label{v4-arith:asi:finite-tensorator}
 \mu_{X,Y}:\mathsf P_X\otimes_{R_0}\mathsf P_Y\longrightarrow
 \mathsf P_{X\otimes Y},\qquad
 f\otimes g\longmapsto a_{X,Y,S_k}^{-1}(1_X\otimes g)f,
\end{equation}
where $f:S_i\to X\otimes S_j$ and $g:S_j\to Y\otimes S_k$.

To describe the retained symmetry, associate to a finite $R_0$-bimodule
$P$ the linear functor
\[
 T_P(S_j)=\bigoplus_i S_i\otimes_{\mathbb C}e_iPe_j.
\]
Here tensoring an object with a vector space means the corresponding
finite direct sum. Extend $T_P$ to sums and maps by linearity.

A right $\mathcal C$-module structure on $T_P$ consists of natural
isomorphisms
\[
 \lambda_{M,Z}:T_P(M\otimes Z)\longrightarrow T_P(M)\otimes Z
\]
compatible with the unit and associator. Explicitly, the two maps from
$T_P((M\otimes Z)\otimes W)$ to
$(T_P(M)\otimes Z)\otimes W$ agree:
\[
 (\lambda_{M,Z}\otimes1_W)\lambda_{M\otimes Z,W}
 =a^{-1}_{T_P(M),Z,W}\lambda_{M,Z\otimes W}T_P(a_{M,Z,W}).
\]
A morphism must commute with every $\lambda_{M,Z}$. These are finite
conditions: simple $M,Z,W$ and bases of their junction spaces determine
all the displayed maps and equations by additivity.

Let $\mathcal S_{\mathcal C}$ be this category. Its tensor product is
composition of module functors, expressed on coefficients by
$P\otimes_{R_0}Q$. For $\mathsf P_X$, the functor is
$L_X(M)=X\otimes M$, and its module structure is
$\lambda^X_{M,Z}=a^{-1}_{X,M,Z}$.

\begin{theorem}\label{v4-arith:asi:finite-coefficients}
The maps \eqref{v4-arith:asi:finite-tensorator} give a tensor equivalence
\[
 \mathcal C\xrightarrow{\ \sim\ }\mathcal S_{\mathcal C},
 \qquad X\longmapsto(\mathsf P_X,\lambda^X).
\]
In particular the symmetry-compatible junctions between these finite
bimodules are exactly the original junctions of $\mathcal C$.
The equivalence preserves the specified duality and spherical evaluations.
\end{theorem}
\begin{proof}
Choose inclusions $g_{j,\alpha}:S_j\to Y\otimes S_k$ and projections
$\bar g_{j,\alpha}$ with
$\sum_{j,\alpha}g_{j,\alpha}\bar g_{j,\alpha}=1$.
The inverse of \eqref{v4-arith:asi:finite-tensorator} sends
$h:S_i\to(X\otimes Y)\otimes S_k$ to
\[
 \sum_{j,\alpha}
 (1_X\otimes\bar g_{j,\alpha})a_{X,Y,S_k}h
       \otimes g_{j,\alpha}.
\]
The two composite identities follow from this resolution of the identity.
They also prove independence of its choice. The associativity coherence
is the pentagon for $a$, and the unit coefficient is $\mathsf P_1=R_0$.

Every linear functor of the finite semisimple category is recovered
from its matrix of spaces $\operatorname{Hom}(S_i,T(S_j))$.
The same resolution proves that composition of functors is the balanced
product of their matrices. Thus the definition of $\mathcal S_{\mathcal C}$
retains exactly the module functors and their module-natural maps.

For a module functor $T$, its structure map at $(1,M)$ identifies
$T(M)$ naturally with $T(1)\otimes M$. The module-functor equation
makes this identification compatible with the right action.
For a module-natural map $\eta:L_X\to L_Y$, the same equation gives
$\eta_M=\eta_1\otimes1_M$, with the unit constraints understood.
Evaluation at $1$ is therefore an inverse equivalence, on both objects
and maps. Composition identifies $L_XL_Y$ with $L_{X\otimes Y}$
through $a^{-1}$. This proves the tensor assertion.

Transport the given evaluations, coevaluations, and pivotal maps
through this equivalence. Their defining equations are preserved,
as is every closed composition. This gives the asserted duality.
\end{proof}

For the near-group rules the size of the construction is explicit.
Order the simples by $G$ followed by $\rho$. The matrix of dimensions
of $\mathsf P_\rho$ is
\begin{equation}\label{v4-arith:asi:finite-near-group}
 D_\rho=\begin{pmatrix}0_{N\times N}&\mathbf1\\
                         \mathbf1^{\mathsf t}&N\end{pmatrix}.
\end{equation}
For $g\in G$, $D_g$ is the permutation $h\mapsto g+h$ on the
$G$-block and is the identity on the $\rho$-entry. Hence
$\dim R_0=N+1$ and $\dim\mathsf P_\rho=3N$.

If $\mathbf d=(1,\ldots,1,d)^{\mathsf t}$, then
\[
 D_g\mathbf d=\mathbf d,\qquad D_\rho\mathbf d=d\mathbf d,
 \qquad d^2=N+Nd.
\]
The last identity also explains the relative dimension. Closing a
coefficient line over the $j$-th boundary weights its $i$-th summand
by $d_i/d_j$. The total is $d_X$. These weights are the regular-module
trace $\operatorname{tr}_{S_i}(1)=d_i$ of the given spherical category.
They need not be positive.

The symmetry condition cannot be suppressed. Ordinary bimodule maps
$R_0\to R_0$ form $\mathbb C^{\mathcal I}$. A module-natural endomorphism of
the identity has scalars $c_i$ satisfying $c_i=c_0$ by evaluating its
compatibility at $(1,S_i)$. Its space is consequently $\mathbb C$,
as required by the tensor unit of $\mathcal C$.

\subsection{The reconstructed weak Hopf action}
\label{v4-arith:asi:weak-hopf}

The categorical symmetry has an equivalent description by a finite
weak Hopf algebra. It uses the left-insertion tensorator
\eqref{v4-arith:asi:finite-tensorator}. The choice matters for its
associator coefficients.

Let $\mathsf P$ denote the vector-space-valued functor underlying the
preceding construction and put
\[
 H_{\mathcal C}=\operatorname{Nat}(\mathsf P,\mathsf P)
       =\bigoplus_{s\in\mathcal I}\operatorname{End}_{\mathbb C}(\mathsf P_{S_s}).
\]
Multiplication is composition. These matrices need not be
\(R_0\)-bilinear. An element \(h\) acts on every
$\mathsf P_X$, with its action on sums and retracts determined by
naturality.

Let $\pi_{X,Y}$ be the balanced quotient and let
\[
 j_{X,Y}(p\otimes_{R_0}q)=\sum_i pe_i\otimes e_iq
\]
be its section. The coproduct is uniquely specified by its action on
all pairs of simple representations:
\begin{equation}\label{v4-arith:asi:weak-coproduct}
 (\Delta h)_{X,Y}
   =j_{X,Y}\mu_{X,Y}^{-1}h_{X\otimes Y}\mu_{X,Y}\pi_{X,Y}.
\end{equation}
In particular $\Delta(1)=E$, where
$E(p\otimes q)=\sum_i pe_i\otimes e_iq$.
With $\epsilon_0(e_i)=1$, its counit is
\[
 \epsilon(h)=\epsilon_0(h_1(1_{R_0})),\qquad
 \epsilon(1)=|\mathcal I|,
\]
where \(\mathcal I\) denotes the chosen set of simple labels.

These formulas define a weak Hopf algebra. Here ``weak'' means that
the coproduct preserves multiplication but its unit is the idempotent
$E$, and the counit obeys the corresponding weak counit identities.
Its representation tensor product is $E(U\otimes V)$, with the
action given by $\Delta$. Its unit representation is $R_0$.
This is the dual form of Hayashi's finite semisimple reconstruction,
Theorems~3.1 and~4.1.

The identities are visible in the formulas. Multiplicativity uses
$\pi j=1$. Coassociativity uses the associativity of $\mu$ and
the equality of the two projections onto matching triples of boundary
idempotents. The weak unit and counit equations reduce to
$e_ie_j=\delta_{ij}e_i$ and $\epsilon_0(e_i)=1$.

The antipode is the transpose action on the dual coefficients.
More explicitly, let
\[
 b_X(q,p)=\epsilon_0\mathsf P(\operatorname{ev}_X)
               \mu_{X^*,X}\pi_{X^*,X}(q\otimes p).
\]
This pairing between $\mathsf P_{X^*}$ and $\mathsf P_X$ is
nondegenerate. Define $S(h)$ by
$b_X(q,S(h)_Xp)=b_X(h_{X^*}q,p)$.
The duality snake equations give the antipode equations, with the
source and target counital projections in place of the scalar unit.

Every finite $H_{\mathcal C}$-module is a sum of the displayed simple
modules $\mathsf P_{S_s}$. Its intertwiners are precisely the maps
coming from $\mathcal C$, and \eqref{v4-arith:asi:weak-coproduct} identifies
its tensor product with \eqref{v4-arith:asi:finite-tensorator}.
Thus $\mathcal C\simeq\operatorname{Rep}(H_{\mathcal C})$ as a
tensor category. Equivalently, one can impose this weak Hopf action
instead of the module-functor symmetry maps.

For the near-group rules,
\[
 \dim H_{\mathcal C}=N(N+1)^2+9N^2.
\]
At $N=2$ this is $54$. On
$\mathsf P_\rho\otimes\mathsf P_\rho$, the idempotent $E$ has rank
$18$ in dimension $36$. Its role cannot be replaced by the identity.

\subsection{The associator in a pointed example}
\label{v4-arith:asi:pointed}

Let \(G_0\) be a finite abelian group and let
\(\Omega:G_0^3\to\mathbb C^\times\) be a normalized three-cocycle.
In \(\mathrm{Vec}_{G_0}^{\Omega}\), the simple label \(g\) has
coefficient basis \(v_{g,j}\), with left endpoint \(g+j\) and
right endpoint \(j\). The tensorator is
\[
 \mu(v_{g,h+j}\otimes_{R_0}v_{h,j})
       =\Omega(g,h,j)^{-1}v_{g+h,j}.
\]
Thus a nontrivial associator survives even though every simple
fusion multiplicity is one.

Write \(E^g_{ab}\) for the matrix sending \(v_{g,b}\) to
\(v_{g,a}\) and killing the other vectors and simple blocks.
The reconstruction gives
\begin{align}
 E^g_{ab}E^h_{cd}
   &=\delta_{g,h}\delta_{b,c}E^g_{ad},\notag\\
 \Delta(E^s_{ab})
   &=\sum_{g+h=s}
      \frac{\Omega(g,h,a)}{\Omega(g,h,b)}
      E^g_{h+a,h+b}\otimes E^h_{ab},
       \label{v4-arith:asi:pointed-coproduct}\\
 \epsilon(E^g_{ab})&=\delta_{g,0},\notag\\
 S(E^s_{ab})
   &=\frac{\Omega(s,-s,b+s)}{\Omega(s,-s,a+s)}
         E^{-s}_{b+s,a+s}.\notag
\end{align}
The counit is nonzero on every matrix unit in the unit block,
including the off-diagonal ones.

Indeed, composing \(\mu\), the matrix unit and \(\mu^{-1}\)
contributes first \(\Omega(g,h,b)^{-1}\), then
\(\Omega(g,h,a)\). Matching the endpoints gives the displayed
matrix indices. Every matrix unit in the unit block sends
\(1_{R_0}\) to a vector of coordinate sum one.
The dual pairing pairs \(v_{-g,j+g}\) with \(v_{g,j}\), with
factor proportional to \(\Omega(-g,g,j)^{-1}\).
Its transpose gives the antipode. The evaluation normalization
cancels in the ratio.

The identity
\[
 \Omega(h,k,j)\Omega(g,h+k,j)\Omega(g,h,k)
 =\Omega(g+h,k,j)\Omega(g,h,k+j)
\]
is the equality of the iterated tensorators. It proves
coassociativity directly. The matching-endpoint projections and
duality snakes give the remaining weak Hopf identities.

For \(G_0=\mathbb Z/n\), choose representatives
\(0,\ldots,n-1\) and a primitive root \(\zeta_n\). The carry cocycle
is
\[
 \Omega(g,h,j)=\zeta_n^{\,g\lfloor(h+j)/n\rfloor}.
\]
The cocycle equation expresses equality of the total carry under
two parenthesizations. At \(n=3\), the \((g,h)=(1,2)\) summand of
\(\Delta(E^0_{01})\) and one antipode value are
\[
 \zeta_3^{-1}E^1_{20}\otimes E^2_{01},
 \qquad
 S(E^1_{02})=\zeta_3^{-1}E^2_{01}.
\]
These values distinguish the ratio in
\eqref{v4-arith:asi:pointed-coproduct} from its inverse.
Even when \(\Omega=1\), \(\Delta(1)\) remains the
matching-endpoint projection.

\section{Finite-dimensional conformal grades for the general family}
\label{v4-arith:asi:sec:finite-analytic}

We now combine this finite coefficient carrier with an explicit
conformal theory. Let $\Lambda=E_8$ be the positive definite lattice
\[
 \Lambda=\{x\in\mathbb Z^8\cup(\mathbb Z+\tfrac12)^8:
                                      \textstyle\sum_i x_i\in2\mathbb Z\}.
\]
Its even unimodular form gives the holomorphic lattice vertex
algebra \(W=V_\Lambda\) of central charge eight, as in Dong--Mason.

Its full space of states is
\[
 W=\bigoplus_{\alpha\in\Lambda}
 \operatorname{Sym}\left(\bigoplus_{n>0}\mathbb C^8t^{-n}\right)e^\alpha.
\]
A monomial $h^1_{-n_1}\cdots h^r_{-n_r}e^\alpha$ has weight
$\sum n_i+(\alpha,\alpha)/2$.
Positive definiteness makes each weight space finite dimensional.

For completeness, choose a normalized lattice cocycle $\epsilon$ with
$\epsilon(\alpha,\beta)/\epsilon(\beta,\alpha)=(-1)^{(\alpha,\beta)}$.
The currents obey
$[h_m,k_n]=m(h,k)\delta_{m+n,0}$, and the lattice fields are
\[
 \begin{aligned}
 Y(e^\alpha,z)
 &=\exp\left(\sum_{n>0}\frac{\alpha_{-n}}n z^n\right)
   \exp\left(-\sum_{n>0}\frac{\alpha_n}n z^{-n}\right)
   e_\alpha z^{\alpha_0},\\
 e_\alpha e^\beta&=\epsilon(\alpha,\beta)e^{\alpha+\beta}.
 \end{aligned}
\]
Normal products with differentiated currents define the fields of all
other states. The translation is $L_{-1}$ and the conformal vector is
$\frac12\sum_{a=1}^8(h^a_{-1})^2 1$ for an orthonormal basis $(h^a)$.

For $(P,\lambda)\in\mathcal S_{\mathcal C}$ define
$\mathsf V(P)=W\otimes P$. Retain the $R_0$ actions and the finite
module-functor maps $\lambda$. A field is required to respect them.
For a compatible junction $f:P\otimes_{R_0}Q\to S$ it is
\begin{equation}\label{v4-arith:asi:finite-analytic-field}
 \mathcal J_f(u\otimes p,z)(v\otimes q)
       =Y_W(u,z)v\otimes f(p\otimes_{R_0}q).
\end{equation}


The algebra of boundary vacua is $W\otimes R_0$. Its weight-zero
space is $R_0$. Completions are products over conformal weights,
and states themselves remain finite sums.

Equivalently, $H_{\mathcal C}$ acts on the coefficient factors and the
fields obey the concrete covariance equation
\[
 h\,\mathcal J_f(u\otimes p,z)(v\otimes q)
 =\sum_{(h)}\mathcal J_f(u\otimes h_{(1)}p,z)
                           (v\otimes h_{(2)}q),
 \qquad \Delta h=\sum_{(h)}h_{(1)}\otimes h_{(2)}.
\]
The weak coproduct projects to the matching boundary sectors before
the field is applied.

\subsection{The sewing domain}
\label{v4-arith:asi:sewing-domain}

A sewing datum starts with finitely many compact pointed Riemann
surfaces, possibly disconnected. At each paired point choose
coordinate discs of radii \(r_j,r'_j\), with all discs pairwise
disjoint. Glue their coordinates by \(u_jv_j=t_j\), where
\[
 0<|t_j|<r_jr'_j.
\]
Each initial connected component must have a marked point that
remains unsewn. Fix the local coordinates, contragredient pairings
and the initial conformal block. Fractional powers of \(t_j\)
use a chosen branch of \(\log t_j\).

For a \(C_2\)-cofinite vertex algebra and finitely generated modules,
Gui's theorem gives absolute, locally uniform
convergence of the full propagator sums on this domain.
Here \(C_2\)-cofiniteness means that the quotient by the span of
\(u_{(-2)}v\) is finite dimensional.
The theorem allows simultaneous sewings that create handles.
It concerns the specified conformal block. It does not choose a
globally normalized scalar partition function.

If an initial component has no remaining insertion, first insert
a vacuum. Apply propagation of vacua, then remove the auxiliary
mark through the propagation isomorphism.
The hypothesis is not removed by simply omitting the point.
All finite simultaneous-sewing assertions below use this domain.

\begin{theorem}\label{v4-arith:asi:finite-analytic-realization}
Every arithmetic fusion category considered here has a fully faithful
analytic ordered realization
$X\mapsto\mathsf V(\mathsf P_X,\lambda^X)$ with finite-dimensional
homogeneous state spaces. Its symmetry-compatible relative fusion is
\[
 \mathsf V(\mathsf P_X)\boxtimes_{W\otimes R_0}
 \mathsf V(\mathsf P_Y)
       =\mathsf V(\mathsf P_{X\otimes Y}),
\]
with tensorator \eqref{v4-arith:asi:finite-tensorator} on coefficients.
Its junctions, duality, and all fusion-tree sewings recover $\mathcal C$.
All descendant sewings on the domain of
Section~\ref{v4-arith:asi:sewing-domain} converge absolutely and locally
uniformly, including the permitted sewings that create handles.
Taking full collision principal parts gives relative interfaces in the
ordered word category, with compatible cut-cobar comparisons in every arity.
\end{theorem}
\begin{proof}
The lattice algebra is simple and its vacuum module has scalar
endomorphisms. A $W$-module map between finite sums of this module
is therefore $1_W\otimes f$. The space of vacuum-to-vacuum
intertwining operators is one dimensional, generated by $Y_W$.
Consequently every balanced field between these finite sums has
the form \eqref{v4-arith:asi:finite-analytic-field}. Its balancing and
symmetry identities say precisely that $f$ is a morphism in
$\mathcal S_{\mathcal C}$ out of $P\otimes_{R_0}Q$.
Theorem~\ref{v4-arith:asi:finite-coefficients} now proves full faithfulness,
the tensor assertion, and preservation of the specified duality.

An iterated product is a finite sum of $W$ conformal blocks multiplied
by compositions of coefficient maps. The vertex-algebra associativity
identifies their analytic continuations. In fusion-tree bases the
coefficient comparison is the original associator of $\mathcal C$.
The module-functor pentagon makes this identity compatible with every
retained symmetry map and every further cut.

Positive definiteness makes the lattice algebra regular by
Dong--Li--Mason, Theorem~3.16. Abe--Buhl--Dong, Theorem~4.5,
then gives rationality and \(C_2\)-cofiniteness. All the state
spaces in the construction are finite sums of its vacuum module.
Its propagator is the sum of the identity tensors of $W_n$, multiplied
by $t^n$. The coefficient contractions use the transported relative
evaluations. They are finite dimensional, so each sewn expression is
a finite linear combination of ordinary $W$ sewings.

Gui's Theorem~13.1 applies to each summand with all the hypotheses of
Section~\ref{v4-arith:asi:sewing-domain}. In particular, each initial
connected component retains an unsewn marked point.
It permits simultaneous sewings and disconnected initial surfaces,
including sewings that create handles. Finite linear combinations preserve
absolute and locally uniform convergence. The sewing identities on
coefficients and states consequently remain valid after these sums.

Finally apply the local-current Cousin construction to the ordinary
lattice chiral algebra and denote its cut cobar by $A_W$.
The proof of Theorem~\ref{v4-arith:asi:all-arity} with scalar coefficients
uses only translation, finite pole orders, and chiral Jacobi, all of
which hold for $W$. Set
\[
 \mathfrak R_W=A_W\otimes R_0,\qquad
 \mathfrak P_X=A_W\otimes\mathsf P_X.
\]

The relative bar comparison \eqref{v4-arith:asi:relative-bar} gives
\[
 \mathfrak P_X\star^{\mathbb L}_{\mathfrak R_W}\mathfrak P_Y
 \simeq A_W\otimes(\mathsf P_X\otimes_{R_0}\mathsf P_Y)
 \xrightarrow{1\otimes\mu_{X,Y}}\mathfrak P_{X\otimes Y}.
\]
There are no higher coefficient Tor groups because $R_0$ is semisimple.


The normal principal part of \eqref{v4-arith:asi:finite-analytic-field}
is the lattice-current face tensored with its coefficient junction.
At a nonconsecutive collision the full ordered coefficient tensor
is retained as a spectator. Thus no interchange of the two interface
labels is needed. Every geometric face commutes with the finite
symmetry maps, and the relative bar comparison commutes with every
cut. These identities prove the geometric assertion on the word carrier.
\end{proof}

For example, take a unit-norm current in $W$ and close the arithmetic
junction of \eqref{v4-arith:asi:example-unit} on its coefficient line.
The resulting two-current amplitude is $-e_1(z-w)^{-2}$.
All its higher oscillator and lattice descendants are present.
At $N=2$, the three sector multiplicities are $3,3,6$. Since
$\dim W_1=248$, their first positive grades have dimensions
$744,744,1488$.

This is a theory with specified categorical symmetry and several
boundary vacua. The underlying ordinary lattice-module category does
not encode $\mathcal C$ by itself. The finite symmetry maps are part
of the field theory, and their compatible junctions encode precisely
$\mathcal C$. In particular this construction makes no assertion that
$\mathcal C$ is the category of local modules of a vertex algebra.

\section{The obstruction at the empty configuration}

The defect of restriction along \eqref{v4-arith:asi:theta} is already visible
at the empty configuration. Set
\[
 \mathcal T_R=\mathfrak R\star^{\mathbb L}_{G_R}\mathfrak R,\qquad
 \epsilon_R:\mathcal T_R\to\mathfrak R.
\]
Associativity of relative tensor products gives, for two
$\mathfrak R$-bimodules $P,Q$,
\[
 \operatorname{Res}P\star^{\mathbb L}_{G_R}\operatorname{Res}Q
 \simeq\operatorname{Res}
 \left(P\star^{\mathbb L}_{\mathfrak R}\mathcal T_R
                \star^{\mathbb L}_{\mathfrak R}Q\right).
\]

The comparison with $\operatorname{Res}(P\star^{\mathbb L}_{\mathfrak R}Q)$
uses $\epsilon_R$. At the empty configuration it is multiplication
$R\otimes_F R\to R$. For $r\notin F1$ its kernel contains the nonzero
element $1\otimes r-r\otimes1$. A linear functional taking $1$ to zero
and $r$ to one proves nonvanishing. The restricted relative unit also has
empty coefficient $R$, whereas the $G_R$-bimodule unit has coefficient $F$.

Changing the particular connected bimodules cannot remove this problem.
\begin{theorem}\label{v4-arith:asi:empty-obstruction}
Let $D$ be an algebra in $\mathcal W_X$ with $D_\emptyset\simeq F$ as a
unital algebra.
Let $\mathcal D$ be a rigid category with the fusion rules
\[
 \alpha_g\alpha_h=\alpha_{g+h},\qquad
 \alpha_g\rho=\rho\alpha_g=\rho,\qquad
 \rho^2=\bigoplus_{g\in G}\alpha_g\oplus\rho^{\oplus N},\quad |G|=N.
\]
For $N>0$, there is no additive strong unital tensor functor
$\mathcal D\to\operatorname{Bimod}_D(\mathcal W_X)$ for the derived
balanced cut tensor product.
\end{theorem}
\begin{proof}
Empty evaluation is strong monoidal for that product. Indeed the empty
component of its two-sided bar is
\[
 P_\emptyset\otimes_F D_\emptyset^{\otimes n}\otimes_F Q_\emptyset,
\]
whose realization is $P_\emptyset\otimes_F^{\mathbb L}Q_\emptyset$.
It takes the unit to $F$. A proposed functor would therefore give a
strong tensor functor to the ordinary $\mathbb Z$-graded derived category
$D(F)$.

Images of rigid objects are dualizable. A coevaluation over a field is
a finite sum of homogeneous tensors. The snake identity expresses every
cohomology class in the span of their components. Thus these images are
perfect complexes.

An invertible object of $D(F)$ is a shift of a one-dimensional vector
space. Since each $\alpha_g$ has finite tensor order, its shift is zero
and its Euler characteristic is one.

If $a$ is the integer Euler
characteristic of the image of $\rho$, additivity and multiplicativity
give $a^2=N+Na$. This is impossible because
\[
 (N+1)^2<N^2+4N<(N+2)^2\qquad(N\ge1).
\]
Thus the dimension polynomial has no integer root. No faithfulness
hypothesis was needed.
\end{proof}

The theorem concerns the stated balanced tensor product. Fusion of rooted
probes defined by collision or nearby-cycle operations is a different
tensor product and is not excluded. Likewise, the coefficient exchange
$c_B$ is not a braiding of $\mathcal C$. It need not even be linear for
the left regular $B$-action: the difference between $c_B(ab,c)$ and
$a\,c_B(b,c)$ is $1\otimes abc-a\otimes bc$. These operations occur on
different carriers.

\section{Finite recoupling and the \texorpdfstring{$E_6$}{E6} connection}

The multiplicity space
$M=\operatorname{Hom}(\rho,\rho^2)=\bigoplus_hFT_h$ carries
\[
 p^gT_h=T_{h-g},\qquad q^gT_h=\chi(g,h)T_h.
\]
Let $P_0$ project to $FT_0$, and let $P_{\mathrm{av}}$ project to the
constant vector by averaging. The rank-one operator
$D=P_0P_{\mathrm{av}}$ obeys
\[
 D\left(\sum_hv_hT_h\right)=N^{-1}\sum_hv_hT_0,\qquad D^2=N^{-1}D.
\]
For $N>1$ it is not idempotent. Nor is it the tensor-unit projection in
$\rho^2$: that projection is $X_0X'_0$, which kills every $T_h$.

The corrected identity of Radchenko--Wheeler, Section~4.2,
acts on \(M\). Its right-hand side is \(-N\beta D\), as in
\eqref{v4-arith:mf:corrected-formula}.

The geometric associativity homotopy has a different type. If
$dh=L-R$ in a differential graded operation complex, every chain map from that complex
to $\operatorname{End}_F(M[0])$ sends $L-R$ to zero. It therefore cannot
send this boundary to the nonzero corrected finite-pentagon difference.
A comparison through a larger channel system would have to specify the
intermediate operations and the compression. The two corrections cannot
be identified by a chain map to this degree-zero space.

Allowing negative degrees does not remove this obstruction if the
comparison must preserve the full finite operator algebra.
\begin{proposition}\label{v4-arith:asi:exact-correction}
Let $E$ be a unital differential graded algebra and
$\iota:\operatorname{End}_F(M)\to Z^0(E)$ a unital algebra map.
If $\iota(D)$ is exact, then $E$ is acyclic.
\end{proposition}
\begin{proof}
Write $E_{ij}$ for the matrix units in the basis $(T_g)_{g\in G}$.
The displayed formula for $D$ gives
\[
 E_{i0}DE_{0i}=N^{-1}E_{ii},\qquad
 N\sum_iE_{i0}DE_{0i}=1.
\]
If $dh=\iota(D)$, the degree $-1$ element
$H=N\sum_i\iota(E_{i0})h\iota(E_{0i})$ satisfies $dH=1$,
since the matrix units have closed images. Left multiplication by $H$
contracts $E$: $d(Hx)+H\,dx=x$. Hence every cohomology group vanishes.
\end{proof}

The matrix units are already generated by the finite Heisenberg pair.
Indeed, Fourier averaging of the $q^g$ gives $P_0$, and translates
$p^{-i}P_0p^j$ give $E_{ij}$. Thus a comparison retaining the Heisenberg
operators and their products cannot make the nonzero correction exact
in a recipient with nonzero cohomology. A geometric comparison must
retain the larger channel system and specify how its compression fails
to preserve either differentials or compositions.

\subsection{Compression and the supported operation complex}
\label{v4-arith:asi:compression}

Compression can produce a nonzero difference from an identity on
a larger space. Suppose
\(\mathsf A\mathsf B=\mathsf C\mathsf D\mathsf E\) there.
Let \(P^2=P\) and \(Q=1-P\). Inserting \(P+Q\) at every
intermediate position gives
\begin{align}
 (P\mathsf AP)(P\mathsf BP)
 &-(P\mathsf CP)(P\mathsf DP)(P\mathsf EP)\notag\\
 &=P\mathsf CQ\mathsf DP\mathsf EP
   +P\mathsf CP\mathsf DQ\mathsf EP\notag\\
 &\quad{}+P\mathsf CQ\mathsf DQ\mathsf EP
   -P\mathsf AQ\mathsf BP.
 \label{v4-arith:asi:discarded-channels}
\end{align}
All four terms record discarded intermediate channels.
To obtain \(-N\beta D\), one must construct the larger supported
maps and the specified projection \(P\), then identify their
displayed sum with that operator.

The supported construction retains more than the total diagonal.
At three positions it includes
\[
 [\,\mathcal V_{2|13}\xrightarrow{-r_{2,(13)}}
                       \mathcal V_{123}\,]
\]
in degrees \(-1,0\), with outer-collision homotopy \(r_{13}\).
At four positions it also retains \(13|24\).
Normal derivatives, the positions of spectator states, the
empty cut unit and the local vacuum are part of the comparison.

There is an obstruction calculation for constructing these
operations themselves. Let \(\mathcal E\) be a differential
graded nonsymmetric operad of admissible supported operations.
Require closure under insertions and impose the translation,
differential-module, cut and vacuum conditions in its definition.
Use suspended operations, let \(b_1\) be the differential, and
write \(\partial=[b_1,-]\).

Suppose degree-one operations \(b_2,\ldots,b_{n-1}\) satisfy all
identities below arity \(n\). Define
\[
 \mathfrak o_n=
 \left(\sum_{i=2}^{n-1}b_i\right)^2_{\mathrm{arity}\ n}
       \in\mathcal E(n)^2.
\]
Then \(\partial\mathfrak o_n=0\), and \(b_n\) exists exactly when
\[
 [\mathfrak o_n]=0\quad\text{in }H^2(\mathcal E(n)).
\]
It must satisfy \(\partial b_n=-\mathfrak o_n\).
With the lower operations fixed, solutions modulo normalized
degree-zero homotopy form an \(H^1(\mathcal E(n))\)-torsor.

Indeed, for \(B=b_1+\cdots+b_{n-1}\), associativity of
composition gives \([B,B^2]=0\).
All lower-arity components of \(B^2\) vanish, leaving
\([b_1,\mathfrak o_n]=0\) in arity \(n\).
Adding \(b_n\) changes this component by \(\partial b_n\).
The difference of two solutions is a degree-one cocycle,
and a normalized homotopy changes it by a boundary.

This is the obstruction to a new operation.
The higher-map problem of
Theorem~\ref{v4-thm:recognition-one-step} instead has obstruction
in degree one and choices in degree zero.
Neither cohomology is determined by the rigidity of the finite
fusion category or by a calculation on position-free states.

The full recoupling space is larger and is explicitly available. For the
quadratic example of \eqref{v4-arith:asi:example-unit},
$\operatorname{Hom}(\rho,\rho^3)$ has dimension six. Its left and right
tree bases are
\[
 L_g=X_g,\quad L_{h,k}=T_hT_k,\qquad
 R_g=\rho(X_g)U_g,\quad R_{h,k}=\rho(T_h)T_k.
\]

Their projections are
$L'_g=X'_g$, $L'_{h,k}=T'_kT'_h$,
$R'_g=U_{-g}\rho(X'_g)$ and $R'_{h,k}=T'_k\rho(T'_h)$.
The change-of-basis matrix is $\mathsf F_{ij}=R'_iL_j$. Its inverse has entries
$L'_iR_j$. In the order $0,1,(0,0),(0,1),(1,0),(1,1)$, put
$r=(1+\sqrt3)/2$, $t=(1+i)/2$, $\bar t=(1-i)/2$ and
$c=(1-i)(\sqrt3+i)/4$. Direct substitution gives
\begin{equation}\label{v4-arith:asi:F6}
 \mathsf F=\begin{pmatrix}
 -r&-r&-r&0&0&-ir\\
 -r&r&0&-ir&-r&0\\
 c&c&r&0&0&t\\
 c&-c&0&t&r&0\\
 c&c&\bar t&0&0&ir\\
 -c&c&0&-ir&-\bar t&0
 \end{pmatrix}.
\end{equation}

Put \(a_F=(2+\sqrt3+i)/4\). In the same order, the inverse is
\begin{equation}\label{v4-arith:asi:F6-inverse}
 \mathsf F^{-1}=
 \begin{pmatrix}
 -r&-r&-a_F&-a_F&-a_F&a_F\\
 -r&r&-a_F&a_F&-a_F&-a_F\\
 1&0&r&0&t&0\\
 0&-i&0&\bar t&0&ir\\
 0&1&0&r&0&-t\\
 -i&0&\bar t&0&-ir&0
 \end{pmatrix}.
\end{equation}
Multiplying the two displayed matrices gives the identity.
Equivalently, this follows by inserting the resolutions
\(\sum_iL_iL'_i=\sum_iR_iR'_i=1\) between their entries.

For every number of inputs the analogous tree bases lie in one
intertwiner space. Their change-of-basis matrices obey the pentagon
because both products compare the same two bases. The relative chiral
functor of Theorem~\ref{v4-arith:asi:relative-realization} transports these
actual maps and their coherence.

There is an independent singular realization at affine $\mathfrak{sl}_2$
level ten. Write $V_j$ for its fusion label of highest weight $j$.
The $E_6$ algebra object $V_0\oplus V_6$ has a module tensor
subcategory of type $(\mathbb Z/2,2)$, as in
Thornton, Example 1.2.

Its free-module images are
\[
 V_0\mapsto1,\quad V_2\mapsto\rho,\quad V_4\mapsto g\oplus\rho,\quad
 V_6\mapsto1\oplus\rho,\quad V_8\mapsto\rho,\quad V_{10}\mapsto g.
\]
These follow from free-module adjunction and the even $SU(2)_{10}$ fusion
rules. Its positive dimension is $1+\sqrt3$. It is the
other embedding of the dimension parameter in
\eqref{v4-arith:asi:normalizations}.

The primary three-input blocks leading to $\rho$ have ranks $3,2,1$,
for total highest weights $2,4,6$. They come from
$V_2^{\otimes3}=V_0\oplus3V_2\oplus2V_4\oplus V_6$ before the free-module
functor.

Fix the spin-one basis $v_0,v_1,v_2$ with
$Hv_j=(2-2j)v_j$, $Ev_1=2v_0$, $Ev_2=v_1$, $Fv_0=v_1$ and
$Fv_1=2v_2$. The remaining raising and lowering actions vanish.
Put $w_{abc}=v_a\otimes v_b\otimes v_c$. Use the ordered bases
\[
\begin{array}{c@{\quad}l}
2:&w_{002}-\frac12w_{011}+w_{020},\\
  &4w_{002}-w_{011}-w_{101}+w_{110},\\
  &w_{002}-\frac12w_{101}+w_{200};\\[2pt]
4:&-w_{001}+w_{010},\ -w_{001}+w_{100};\\[2pt]
6:&w_{000}.
\end{array}
\]

On these covariant highest-weight spaces the Knizhnik--Zamolodchikov
(KZ) connection is
\[
 \nabla=d-\frac1{12}\left(\frac{\Omega_{12}}z+
                     \frac{\Omega_{23}}{z-1}\right)dz,\qquad
 \Omega=\tfrac12H\otimes H+E\otimes F+F\otimes E.
\]
For highest weight two one obtains the matrices
\[
 \Omega_{12}=\begin{pmatrix}-2&4&0\\1&0&1\\0&4&-2\end{pmatrix},\qquad
 \Omega_{23}=\begin{pmatrix}-4&-4&0\\0&-2&-1\\0&0&2\end{pmatrix}.
\]
For highest weight four they are
$\left(\begin{smallmatrix}0&2\\2&0\end{smallmatrix}\right)$ and
$\left(\begin{smallmatrix}-2&-2\\0&2\end{smallmatrix}\right)$.
For highest weight six both equal $(2)$.

These matrices are obtained by restricting the ordinary spin-one Casimirs
to the kernel of the total raising operator. The exponents at zero are
$-1/3,-1/6,1/6$, then $-1/6,1/6$, then $1/6$.
The system of conformal-block functionals uses the dual connection.

This differential system has nontrivial local monodromy. Its comparison
with arithmetic fusion requires the module junctions as well as the
ambient affine channels. We construct those junctions first.

\section{The normalized \texorpdfstring{$E_6$}{E6} tensor comparison}
\label{v4-arith:asi:sec:e6-tensor}

Put $s=\sqrt3>0$, $q=e^{\pi i/12}$, and
$[n]=\sin(n\pi/12)/\sin(\pi/12)$. Write $[n]!=\prod_{j=1}^n[j]$ and
$[0]!=1$. Let $\mathcal D_q$ be the semisimplified category of tilting
$U_q(\mathfrak{sl}_2)$-modules, with simple objects $V_0,\ldots,V_{10}$.
The admissible channels of $V_a\otimes V_b$ are
\begin{equation}\label{v4-arith:asi:su2-fusion}
 |a-b|,\ |a-b|+2,\ldots,\min(a+b,20-a-b).
\end{equation}
All recouplings below have even external labels.

We specify the trivalent normalization explicitly. For an admissible
triple put
\[
 D(a,b,c)=
 \frac{[(-a+b+c)/2]![(a-b+c)/2]![(a+b-c)/2]!}
      {[(a+b+c)/2+1]!}.
\]

For $a,b,c,d$ and intermediate labels $e,f$, put
\[
 \begin{split}
 (A_1,A_2,A_3,A_4)
   &=\tfrac12(a+b+e,\ a+d+f,\ c+b+f,\ c+d+e),\\
 (B_1,B_2,B_3)
   &=\tfrac12(a+b+c+d,\ a+c+e+f,\ b+d+e+f).
 \end{split}
\]
The rational Racah normalization is
\begin{equation}\label{v4-arith:asi:q-racah}
 \begin{aligned}
 (F^{abc}_d)_{fe}
 &=(-1)^{(a+b+c+d)/2}[f+1]D(b,c,f)D(a,f,d)\\
 &\quad{}\times
 \sum_{z=\max A_i}^{\min(B_1,B_2,B_3,10)}
 \frac{(-1)^z[z+1]!}{\prod_i[z-A_i]!\prod_j[B_j-z]!}.
 \end{aligned}
\end{equation}
An inadmissible entry is zero. Columns are left-tree inclusions and rows
are right-tree projections. Equivalently, the column of right-tree
projections is $F^{abc}_d$ times the column of left-tree projections.
The counterclockwise exchange in channel $c$ is
\[
 R^{ab}_c=(-1)^{(a+b-c)/2}
 q^{(c(c+2)-a(a+2)-b(b+2))/4}.
\]

Let $A_E=V_0\oplus V_6$. Its unit is the inclusion of $V_0$.
In the normalization \eqref{v4-arith:asi:q-racah}, define its multiplication
by the unit constraints and
\begin{equation}\label{v4-arith:asi:e6-multiplication}
 m_{66}^{\,0}=-\frac13+\frac{s}{6},\qquad
 m_{66}^{\,6}=1.
\end{equation}
Every other component is zero. This is the connected commutative
separable $E_6$ algebra of Kirillov--Ostrik, \S6.

To see the normalization,
write $m_{66}^{\,0}=\alpha$. The channel $e=2$ of the associativity
equation on $V_6^{\otimes3}\to V_6$ gives
\[
 \alpha=-\frac{(F^{666}_6)_{6,2}}{(F^{666}_6)_{0,2}}
        =-\frac13+\frac{s}{6}.
\]
Substitution in the remaining channels gives associativity. The two
nonzero channels of $m_{66}$ have $R^{66}_0=R^{66}_6=1$.
The uniqueness of the $E_6$ algebra then identifies this multiplication
with the conformal-extension algebra.

Write $\mathcal D_q^+$ for the full even-label subcategory and
$\mathcal M_E=(\mathcal D_q^+)_{A_E}$ for its right $A_E$-modules.
Their tensor product is over $A_E$. The left action on a right module
is obtained using the displayed braiding. The free-module functor is
$F_E(V)=V\otimes A_E$. Its three simple images are
\[
 M_1=F_E(V_0),\qquad M_g=F_E(V_{10}),\qquad
 M_\rho=F_E(V_2).
\]


Adjunction gives their endomorphism spaces and the free-module
decompositions displayed above. Every module is a summand of a free
module because $A_E$ is separable. These decompositions therefore
exhaust the simple objects. They give
$M_\rho\otimes_{A_E}M_\rho=M_1\oplus M_g\oplus2M_\rho$.

The free-module junctions have a concrete basis. For
$a,b,c\in\{0,10,2\}$, use adjunction to write
\[
 \operatorname{Hom}_{A_E}(F_E(V_a)\otimes_{A_E}F_E(V_b),F_E(V_c))
 =\operatorname{Hom}_{\mathcal D_q}(V_a\otimes V_b,V_c\otimes A_E).
\]
Its basis element $p^{c;e,t}_{ab}$ projects to $V_e$ and includes it
in $V_c\otimes V_t$, where $t\in\{0,6\}$ and
$e$ occurs in both tensor products. Order these elements by $t$ and
then $e$. In particular the two $\rho\rho\to\rho$ junctions have
$(e,t)=(2,0),(4,6)$.

To compose such maps, move the first $A_E$ factor
past the next input by $R$ and multiply the two $A_E$ factors by
\eqref{v4-arith:asi:e6-multiplication}. This prescription and
\eqref{v4-arith:asi:q-racah} determine every module recoupling.

Let $\sigma:F\to\mathbb C$ send $\sqrt3$ to $-\sqrt3$ and fix $i$.
For the order-two example $F=\mathbb Q(\sqrt3,i)$. Write
$\mathcal C^\sigma=\mathcal C\otimes_{F,\sigma}\mathbb C$.
The arithmetic projection bases are those supplied by
$1,U_g,X'_g,T'_h$ in the Leavitt presentation. Unit junctions have
their unit normalization on both sides.

We give the comparison on the six remaining binary spaces. Put
\[
 a_* =\frac{(1-i)(-s+i)}2,\qquad
 S=\begin{pmatrix}1&1\\1&-1\end{pmatrix},
\]
and
\[
 P=\begin{pmatrix}
 1&\dfrac{(1-i)(3+s(2-i))}{2}\\[3pt]
 -\dfrac{(1+i)(s+1)}2&\dfrac{(1-i)(3+si)}2
 \end{pmatrix}.
\]
Define changes of projection basis $N^{\mathrm a},N^{\mathrm q}$ by
\[
\begin{array}{c|c|c}
 (a,b;c)&N^{\mathrm a}_{ab;c}&N^{\mathrm q}_{ab;c}\\ \hline
 (g,g;1)&1&-\frac12+\frac{s}{3}\\
 (g,\rho;\rho)&1&-i\\
 (\rho,g;\rho)&1&1\\
 (\rho,\rho;1)&a_*&1-\frac{s}{2}+\frac i2\\
 (\rho,\rho;g)&a_*&-3-3i(2+s)\\
 (\rho,\rho;\rho)&S&P
\end{array}
\]
and set
\begin{equation}\label{v4-arith:asi:e6-junction-map}
 K_{ab;c}=(N^{\mathrm a}_{ab;c})^{-1}N^{\mathrm q}_{ab;c},
 \qquad p^{\mathrm a}_{ab;c}\longmapsto K_{ab;c}p^{\mathrm q}_{ab;c}.
\end{equation}
The projection bases in this equation are columns. The inclusions
transform by the inverse matrices, preserving their pairing.

\begin{theorem}\label{v4-arith:asi:e6-tensor-equivalence}
The junction maps \eqref{v4-arith:asi:e6-junction-map} extend to a tensor
equivalence
\[
 \mathcal E:\mathcal C^\sigma\xrightarrow{\ \sim\ }\mathcal M_E,
 \qquad (1,g,\rho)\longmapsto(M_1,M_g,M_\rho).
\]
They preserve the spherical duality and every marked network value.
\end{theorem}

\begin{proof}
Fix a triple $(a,b,c)$ and output $d$. Applying $K$ at both vertices
changes the projection bases by matrices $L_K$ and $R_K$ for the two trees.
The required equation is
\begin{equation}\label{v4-arith:asi:all-f-comparison}
 F^{\mathrm a}L_K=R_KF^{\mathrm q}.
\end{equation}


The quantum computation takes place in
\[
 \operatorname{Hom}_{\mathcal D_q}
       ((V_a\otimes V_b)\otimes V_c,V_d\otimes A_E).
\]
Choose its ambient basis by the output label \(v\in\{0,6\}\),
the total channel \(j\), and the first-pair channel \(x\).
For a left module tree, write its two junction labels as
\((x,t)\) and \((y,u)\), with intermediate simple \(V_e\).
The coefficient vanishes unless the ambient first-pair channel
is the first junction channel \(x\). Its remaining coefficient is
\[
 \left(\sum_f
    (F^{ect}_j)^{-1}_{yf}\,R^{tc}_f\,(F^{etc}_j)_{fx}\right)
       (F^{dut}_j)_{vy}\,m_{ut}^{\,v}.
\]

For a right module tree with junction labels \((f,t)\),
\((y,u)\), the coefficient is
\[
 (F^{abc}_j)_{fx}\,(F^{aet}_j)^{-1}_{yf}\,
       (F^{dut}_j)_{vy}\,m_{ut}^{\,v}.
\]
Inadmissible channels contribute zero. These are respectively
the crossing of the first algebra factor and the two
associativity changes before multiplying the algebra factors.
Writing their columns as \(L\) and \(R\) gives
\(F^{\mathrm q}=(L^{-1}R)^{\mathsf t}\).

Apply the two normalizations $N^{\mathrm a}$ and $N^{\mathrm q}$ to the Leavitt
composites and to the preceding free-module prescription. Both give
the following matrices. Every unit block and every block not listed
is the identity. Zero-dimensional blocks are omitted:
\[
\begin{gathered}
 F^{g\rho g}_\rho=F^{\rho g\rho}_g=-1,\qquad
 F^{g\rho\rho}_\rho=\begin{pmatrix}0&1\\1&0\end{pmatrix},\\
 F^{\rho g\rho}_\rho=\begin{pmatrix}1&0\\0&-1\end{pmatrix},\qquad
 F^{\rho\rho g}_\rho=\begin{pmatrix}0&-i\\i&0\end{pmatrix},\\
 F^{\rho\rho\rho}_1=u\begin{pmatrix}1&b\\1&-b\end{pmatrix},\qquad
 F^{\rho\rho\rho}_g=u\begin{pmatrix}b&1\\-b&1\end{pmatrix},\\
 b=-i,\qquad u=\frac{(1-i)(1-si)}4.
\end{gathered}
\]


For the remaining block put
\[
 \varphi=\frac{s-1}{2},\quad
 w=\frac{(s-1)(1+is)}8,\quad
 y=\frac{1+i(s-2)}4,\quad z=\frac{-s-i}{4}.
\]
In the channel order used in \eqref{v4-arith:asi:F6}, that block is
\begin{equation}\label{v4-arith:asi:e6-common-f}
 F^{\rho\rho\rho}_\rho=
 \begin{pmatrix}
 \varphi&\varphi&-wb&w&w&-wb\\
 \varphi&-\varphi&-wb&w&-w&wb\\
 1&1&-yb&z&y&-zb\\
 1&1&-zb&y&z&-yb\\
 1&-1&-yb&z&-y&zb\\
 -1&1&zb&-y&z&-yb
 \end{pmatrix}.
\end{equation}


For example the arithmetic block is obtained from
\eqref{v4-arith:asi:F6} by $\sqrt3\mapsto-s$, followed by the two-vertex
changes induced by $a_*$ and $S$. The quantum block is obtained by
the free-module composition with $m_{66}^0=-1/3+s/6$. The displayed
scalar, two-dimensional, and six-dimensional blocks exhaust the
possible triples. They prove \eqref{v4-arith:asi:all-f-comparison}.

All the matrices $K$ are invertible. Equation
\eqref{v4-arith:asi:all-f-comparison}, together with the identity on unit
junctions, is the associativity and unit equation for a tensor functor
between the two semisimple categories. It therefore defines
$\mathcal E$ on sums and matrices of simple objects. The three simple
images are distinct and exhaustive, so the functor is an equivalence.

There is no nontrivial tensor automorphism of the identity in either
category. Indeed its scalar on $\rho$ satisfies both
$\lambda_\rho^2=1$ and $\lambda_\rho^2=\lambda_\rho$, and
$\lambda_g\lambda_\rho=\lambda_\rho$. Thus all scalars are one.


The spherical structure is consequently unique. Its dimension on
$\rho$ is $\sigma(d)=1+s$, equal to
$\dim_{\mathcal D_q}(M_\rho)/\dim_{\mathcal D_q}(A_E)$.
Transporting duality along the tensor equivalence preserves this
structure. In particular it preserves cups, caps, and their closed
composites.
\end{proof}

The embedding $\sigma$ is essential. Before changing the embedding,
the global dimension is $2+(1-s)^2=6-2s$. That of $\mathcal M_E$ is
$6+2s$. These are invariants of the fusion categories, so there is no
complex-linear tensor equivalence at the original embedding. Replacing
$i$ by $-i$ as well also changes the category: the trace of
$F^{\rho\rho\rho}_1$ changes from $(1-is)/2$ to $(1+is)/2$.
It cannot be removed by a binary-junction change of basis.

\section{Affine descendants and analytic sewing}
\label{v4-arith:asi:sec:affine-sewing}

Let $H_j=L_{\widehat{\mathfrak{sl}}_2}(10,j)$, $0\leq j\leq10$.
Its lowest-weight space is the ordinary $\mathfrak{sl}_2$-module of
highest weight $j$, and its conformal weights are
\[
 h_j+n,\qquad h_j=\frac{j(j+2)}{48},\quad n\geq0.
\]
The states are finite sums of negative current modes applied to that
space, modulo the integrable null relation
\begin{equation}\label{v4-arith:asi:affine-null}
 E_{-1}^{\,11-j}v_j=0
\end{equation}
and the submodule it generates. All homogeneous spaces are finite
dimensional. Outputs lie in their product completion. The translation
is $L_{-1}$, and the Sugawara central charge is $5/2$.

Huang--Lepowsky, Theorems~3.7--3.13, give these modules their
vertex tensor category \(\mathcal D_{10}\).

Fix the ribbon equivalence
\[
 \Phi:\mathcal D_q\xrightarrow{\ \sim\ }\mathcal D_{10},
 \qquad \Phi(V_j)=H_j,
\]
with counterclockwise braiding and vacuum unit normalization.
We use Finkelberg's equivalence with the correction in his erratum.
The corrected equivalence uses rigidity and the prescribed opposite
tensor-category convention before fixing the counterclockwise
ribbon structure. Its numerical range includes
\(\kappa=k+h^\vee\ge h^\vee+3\).
For \(\mathfrak{sl}_2\) at \(k=10\), one has
\(h^\vee=2\) and \(\kappa=12\), so the bound holds.

The rational trivalent basis \eqref{v4-arith:asi:q-racah} is transported
through this equivalence. This specifies the normalization of every
affine intertwining operator used below. A basis obtained by separately
setting all three-point highest-weight coefficients to one need not
have that normalization.

Write $\Phi_2:\Phi(X)\boxtimes\Phi(Y)\to\Phi(X\otimes Y)$ for its
tensor comparison. The algebra $\widehat A_E=\Phi(A_E)$ is the conformal
extension $L_{\widehat{\mathfrak{sp}}_4}(1,0)$. Its underlying
$H_0$-module is $H_0\oplus H_6$ (Kirillov--Ostrik, Example~5.3 and Section~6).


For an $A_E$-module $M$, let $U(M)$ denote its underlying object of
$\mathcal D_q$ and set $\widehat M=\Phi(U(M))$. In particular,
\begin{equation}\label{v4-arith:asi:e6-full-states}
 \begin{split}
 \widehat M_1&=H_0\oplus H_6,\\
 \widehat M_g&=H_4\oplus H_{10},\\
 \widehat M_\rho&=H_2\oplus H_4\oplus H_6\oplus H_8.
 \end{split}
\end{equation}
These are full affine state spaces. The last is an ordered, nonlocal
$\widehat A_E$-sector. It is not an ordinary local module of the
extension.

We next define its fields and its relative sewing, including the
balancing map. Choose the special Frobenius coproduct of $A_E$ with
$m\Delta=1_{A_E}$, and put $e=\Delta\eta$. It is uniquely determined
by $m$ and by the counit $(3+s)$ times projection to $V_0$.
Explicitly, put $D_E=3+s$ and $\alpha=-1/3+s/6$. In the inclusion
bases dual to the projection bases its nonzero components are
\[
 \Delta_0^{00}=\Delta_6^{06}=\Delta_6^{60}=D_E^{-1},\qquad
 \Delta_0^{66}=\frac{D_E-1}{D_E\alpha},\qquad
 \Delta_6^{66}=\frac{D_E-2}{D_E}.
\]


For right modules $M,N$ let $r_M,r_N$ be their actions and
$\ell_N=r_Nc_{A_E,N}$. The balancing idempotent is
\begin{equation}\label{v4-arith:asi:balancing-idempotent}
 P_{M,N}=(r_M\otimes\ell_N)
          (1_M\otimes e\otimes1_N).
\end{equation}
Its image is $M\otimes_{A_E}N$. Denote the quotient onto this image
by $q_{M,N}$.

For a junction $p:M\otimes_{A_E}N\to P$, define
\begin{equation}\label{v4-arith:asi:e6-full-field}
 \mathcal Y_p=\mathcal Y_{\,\Phi(U(p)q_{M,N})\,\Phi_2}.
\end{equation}
Here $\mathcal Y_f$ is the intertwining operator corresponding to the
$H_0$-module map $f:\widehat M\boxtimes\widehat N\to\widehat P$.
Thus \eqref{v4-arith:asi:e6-full-field} is a definition on all descendants,
through the universal property of the vertex tensor product.
For an arithmetic projection $p^{\mathrm a}_{ab;c,\gamma}$ its operator is
\begin{equation}\label{v4-arith:asi:arith-affine-field}
 \mathcal Y^{\mathrm a}_{ab;c,\gamma}
   =\sum_\beta(K_{ab;c})_{\gamma\beta}
                       \mathcal Y^{\mathrm q}_{ab;c,\beta}.
\end{equation}
The arithmetic sector corresponding to $a$ has the state space
$\widehat M_a$ of \eqref{v4-arith:asi:e6-full-states}.

\begin{theorem}\label{v4-arith:asi:e6-analytic-comparison}
Equations \eqref{v4-arith:asi:e6-full-states}--\eqref{v4-arith:asi:arith-affine-field}
give an analytic ordered realization of $\mathcal C^\sigma$.
The comparison is the identity on the full state spaces and is $K$ on
binary junctions. It preserves translation, conformal grading, relative
duality, collision expansions, and sewing in every arity.
All products and iterates, paired with restricted-dual states,
converge absolutely on their nested-disc domains. Their continuations have the level-ten conformal-block
monodromy. Sewing on the domain of
Section~\ref{v4-arith:asi:sewing-domain} converges absolutely and
locally uniformly, including the allowed sewings that create handles.
\end{theorem}

\begin{proof}
The Frobenius and separability equations make
\eqref{v4-arith:asi:balancing-idempotent} idempotent and identify its image
with the balanced tensor product. They also identify the two iterated
balancing quotients of $M\otimes N\otimes P$. Applying the strong
tensor functor $\Phi$ gives the corresponding identity of affine module
maps. The vertex-tensor associativity isomorphism identifies those
maps with the products and iterates of
\eqref{v4-arith:asi:e6-full-field}. Its defining identity holds for every
state, by Huang--Lepowsky, Theorems 3.7--3.8.

Apply \eqref{v4-arith:asi:all-f-comparison} to that identity. The changes
on its two binary vertices are exactly $L_K$ and $R_K$, so
\eqref{v4-arith:asi:arith-affine-field} identifies the two complete
operator products. Any two binary trees are joined by associativity
moves. The pentagon identifies the resulting maps between trees,
and the unit constraints handle vacuum insertions. This proves the
comparison in every arity. All maps are affine module maps and
preserve $L_{-1}$, $L_0$, and the null submodules.

For clarity, the sewing along a module $H_j$ uses the tensor
\begin{equation}\label{v4-arith:asi:affine-propagator}
 \mathcal P_j(t)=\sum_{n\geq0}t^{h_j+n}
       \sum_\alpha u_{j,n,\alpha}\otimes u^{\vee}_{j,n,\alpha}.
\end{equation}
The second basis is dual to the first in the contragredient module,
with the specified opposite local coordinate. A branch of $\log t$
fixes the fractional powers. Direct sums, the balancing idempotents,
and the relative evaluations specify the analogous sewing for the
three sectors.

No descendant degree is omitted. The identity tensor
in each finite-dimensional grade is independent of its basis.
Thus the junction change $K$ and its dual inverse cancel at an internal edge.
The spherical comparison of
Theorem~\ref{v4-arith:asi:e6-tensor-equivalence} handles cups and caps.

The integrable affine vertex algebra $H_0$ is rational and
$C_2$-cofinite. Its finite sums $\widehat M_a$ are finitely generated
modules. Genus-zero absolute convergence is part of the cited vertex
tensor theorem.

Sew local coordinates $uv=t$ on disjoint
coordinate discs of radii $r$ and $r'$.
The full sums \eqref{v4-arith:asi:affine-propagator} converge
absolutely and locally uniformly by Gui's Theorem~13.1.
Use \(0<|t|<rr'\) and all the hypotheses of
Section~\ref{v4-arith:asi:sewing-domain}. It allows simultaneous
sewings and disconnected initial surfaces. The balancing and junction
maps form finite sums of such conformal blocks, so the same convergence
holds here. Equality before sewing consequently remains equality after
each further sewing.
\end{proof}

This construction also has ordinary boundary fields. For boundary
conditions $M,N\in\{M_1,M_g,M_\rho\}$ their state space is
\[
 \mathsf H_{MN}=\bigoplus_{j=0}^{10}
 \operatorname{Hom}_{A_E}(V_j\otimes M,N)\otimes H_j.
\]
Composition of the module maps gives their operator-product
coefficients, and the corresponding affine intertwining operators give
their position dependence.

If $P=M_1\oplus M_g\oplus M_\rho$, these
spaces are the matrix entries of the internal endomorphism algebra of
$P$ in $\mathcal D_{10}$. Its associativity gives a grading-restricted
conformal open-string vertex algebra by Huang--Kong, Theorem 4.3.
The three boundary identity states account for its three-dimensional
weight-zero space. Its conformal transformations and disc sewings are
those of Huang--Kong, Theorem 5.4.

\section{Collision coefficients of the affine realization}
\label{v4-arith:asi:sec:e6-collisions}

The fractional powers of the affine fields belong to their coefficient
connection. They must be canceled before taking a normal principal
part. For labels $i,j,k$, put $a_{ij}^k=h_k-h_i-h_j$. On this binary
channel write
\[
 \mathcal Y^k_{ij}(u)=u^{a_{ij}^k}\Theta^k_{ij}(u).
\]
On homogeneous input states, $\Theta^k_{ij}$ has integral powers and
a finite lower bound.

The coefficient symbol $\xi_{ij}^k=u^{-a_{ij}^k}$
satisfies $\partial_u\xi_{ij}^k=-a_{ij}^ku^{-1}\xi_{ij}^k$.
The collision map is
\begin{equation}\label{v4-arith:asi:affine-residue}
 r_{ij}(f\xi_{ij}^k;v,w)
 =\operatorname{pp}_{u=0}\bigl(f\Theta^k_{ij}(v,u)w\bigr).
\end{equation}
The derivative of $\xi_{ij}^k$ and the translation identity of the field
make this a map of right differential modules. In particular the term
$a_{ij}^ku^{-1}\Theta^k_{ij}$ in the field derivative is retained.

For more points, the symbols form the inverse conformal-block
coefficient system. It is the finite-rank integrable subbundle of the
Knizhnik--Zamolodchikov connection. One algebraic description is useful.
Let $U_j\subset H_j$ be the lowest-weight space and fix output $l$.
The coefficient at positions $z_1,\ldots,z_r$ is the subspace of
$\operatorname{Hom}_{\mathfrak{sl}_2}(U_{j_1}\otimes\cdots\otimes U_{j_r},U_l)$
whose invariant functional $\tau$ obeys
\begin{equation}\label{v4-arith:asi:block-null}
 \tau\left(\left(\sum_i z_iE^{(i)}\right)^{11-l}b
                \otimes v_l^\vee\right)=0
 \quad\text{for every lowest-space tensor }b.
\end{equation}
The vector $v_l^\vee$ is highest in the module at infinity.


This is the null relation \eqref{v4-arith:asi:affine-null} moved from infinity
by the current Ward identity. The integrable null relations at the
other points give the same kernel. It is therefore an algebraic bundle with its induced flat
connection, by Beauville, Sections~3--4.

Let $\mathcal K^l_{\mathbf j}$ denote this coefficient bundle, with the
inverse-field convention just specified. On a stratum with blocks
$\pi$, the coefficient with states is
\begin{equation}\label{v4-arith:asi:affine-cousin-term}
 (i_\pi)_*(j_\pi)_*\bigl(
 \bigoplus_{\mathbf j}\mathcal K^l_{\mathbf j}\otimes
       \boxtimes_{B\in\pi}(\omega_X\otimes H_{j_B})\bigr)[|\pi|].
\end{equation}
Sewing frames decompose this bundle near a collision over the allowed
intermediate labels. The faces in these frames are
\eqref{v4-arith:asi:affine-residue}, with the determinant signs used in
\eqref{v4-arith:asi:coefficient-complex}.

\subsection{Algebraic coefficients and radial sewing frames}
\label{v4-arith:asi:algebraic-cuts}

The bundle in \eqref{v4-arith:asi:block-null} is algebraic.
Its defining matrix is polynomial in the positions, and its kernel
has the constant rank given by the integrable fusion multiplicity.
A matrix of constant rank on the reduced configuration variety has
locally free kernel, image and cokernel. Thus the kernel is locally
a direct summand of the constant lowest-state tensor bundle.

It is invariant under common translation and nonzero dilation of
the positions. Under translation by \(c\), the operator
\(\sum_i z_iE^{(i)}\) becomes
\(\sum_i z_iE^{(i)}+c\sum_iE^{(i)}\).
The summands commute. Invariance moves each occurrence of the total
raising operator to the highest vector at infinity, where it
vanishes. Dilation multiplies the null equation by a nonzero scalar.

The Casimir connection preserves this algebraic kernel.
Locally, analytic horizontal conformal blocks span every fiber.
Writing an algebraic section in such a frame shows that its
covariant derivative lies in the same kernel at every point.
It is therefore an algebraic one-form with values in the subbundle.
The inherited connection has regular singularities along diagonals.

Consider a cut with left and right roots \(y',y''\), and put
\(d=y'-y''\ne0\). Contract the two groups by
\[
 y_B=y'+\lambda v_B\quad(B\text{ on the left}),\qquad
 y_B=y''+\lambda v_B\quad(B\text{ on the right}).
\]
The root coordinates have \(v_B=0\). The \(v_B\) record the
internal shapes and \(\lambda\) their common scale.
Let \(A_c\) be the ring of rational functions of these shapes
and roots, with poles allowed only on internal diagonals and \(d=0\).

For total output \(l\), the fused coefficient is
\[
 \mathcal K_c^l
 =\bigoplus_{a,b=0}^{10}
   \mathcal K^l_{ab}\otimes
   \mathcal K^a_{\mathbf j'}\otimes
   \mathcal K^b_{\mathbf j''}.
\]
Only the integrable channels occur. Compose the integrable germs of
the two groups with the two-root germ. The result gives a radial
sewing map
\begin{equation}\label{v4-arith:asi:radial-frame}
 J_c:\mathcal K_c^l\otimes A_c[[\lambda]]
       \longrightarrow
       \widehat{\mathcal K}^l_{\mathbf j'\mathbf j''}.
\end{equation}
Its constant term is the composition of the three lowest-state
maps. Ward recursion on the integrable modules defines every
higher coefficient, including at resonant levels.

\begin{proposition}\label{v4-arith:asi:radial-sewing}
The map \eqref{v4-arith:asi:radial-frame} is an isomorphism over
\(A_c[[\lambda]]\), intertwining the fused and full connections.
Frames for nested collisions and cuts compose.
For fixed shape and roots it converges when
\[
 |\lambda|\max_B|v_B|<|d|/4.
\]
It continues to a holomorphic isomorphism over the region where
every left position has larger real part than every right position.
\end{proposition}

\begin{proof}
Insertion of the three germs gives rational Taylor coefficients,
with precisely the allowed poles. Ward identities hold at every
original point. Uniqueness of the integrable germ determines the
series coefficient by coefficient.

The constant term is injective by semisimple composition of
lowest-state maps. Locally choose a retraction \(L_0\) of this
constant term. Then \(L_0J_c=1+O(\lambda)\) is invertible, and
\((L_0J_c)^{-1}L_0\) is a left inverse.
The image lies in the kernel of the polynomial null-equation
matrix.

Its rank equals the rank of that kernel, since
\[
 \sum_{a,b}N^l_{ab}
             N^a_{\mathbf j'}N^b_{\mathbf j''}
   =N^l_{\mathbf j'\mathbf j''}
\]
by associativity of integrable fusion.
The quotient of the ambient free module by the split image is
projective. The null kernel modulo the image is a rank-zero
submodule of this torsion-free quotient, so it vanishes.
This proves the formal isomorphism.

Two nested frames give the same inserted Ward germ.
Their coefficients are finite sums in each state degree, so
associativity of insertion identifies the series.
Translation of the germ identifies the connections.

Contracting horizontality with \(\lambda\partial_\lambda\)
gives a regular-singular linear equation
\[
 \lambda\frac{dJ_c}{d\lambda}
   =J_c\bigl(C+R(\lambda)\bigr)-C^0J_c.
\]
Here \(C\) is the sum of internal Casimirs and \(C^0\) its fused
action. Each cross term of \(R\) has the form
\[
 t_{ij}\frac{\lambda(v_i-v_j)}
                  {d+\lambda(v_i-v_j)},\qquad
 t_{ij}=\frac{\Omega_{ij}}{12}.
\]
It is holomorphic for
\(|\lambda|\max|v_B|<|d|/2\).
Above the finitely many differences of eigenvalues of \(C\)
and \(C^0\), the recursion has an inverse bounded by a constant
times the inverse degree. The usual power-series majorant
therefore converges on every smaller disc.

The finitely many resonant degrees do not affect convergence:
their coefficients have already been fixed by the germs.

Scale contraction keeps the two groups separated in real part.
Horizontal continuation from the convergent neighborhood extends
\(J_c\) through that region. Flatness and uniqueness identify
the continuations on intersecting neighborhoods.
\end{proof}

Resonance is a reason to use the germs. For example, at level four
the three spin-one inputs with spin-one output have pair exponents
\(-2/3,-1/3,1/3\). The degree-one recursion between the first and
third channels is singular. Its right side vanishes, since one
Lie-algebra generator changes the intermediate spin by at most one.
The differential equation alone leaves a free coefficient.
The specified integrable germ selects it.

\subsection{The coupled product on full cut opens}

For this construction a state system records a color in
\(\{0,\ldots,10\}\) and a state factor on every block of every
partition. Its coefficient diagram attaches
\(\mathcal K^l_{\mathbf j}\) to the entire ordered color word.
It retains the state factors, densities and block-number shift
of \eqref{v4-arith:asi:affine-cousin-term}.
The maps are generated by maps on states and by collision residues
with arbitrary spectator blocks.

Concatenate the two color words and their state factors.
Attaching the coefficient bundle of this whole word defines
the coupled cut product:
\begin{equation}\label{v4-arith:asi:coupled-product}
 \mathcal K\mathscr S'\star_{\mathcal K}\mathcal K\mathscr S''
     =\mathcal K(\mathscr S'\boxtimes\mathscr S'').
\end{equation}
The state tensor is external, but the coefficient bundle includes
all cross-Casimir terms. The unit is \(\mathbb C\) in color zero
at the empty word.

A residue of one factor extends by the residue in the enlarged
configuration, with the other states as spectators.
The radial comparison determines this extension.
Rational expansion is injective, state modules are free,
and coefficient bundles lie inside constant bundles.
The nested-frame identity proves compatibility with compositions
of the generating maps.

Thus \(\star_{\mathcal K}\) is a graded monoidal product. The Koszul rule makes its product differential
square to zero. Both parenthesizations carry the same
whole-word bundle, so the associator on these diagrams is the identity.

For a consecutive cut \(c\), project onto the partitions having no
block that crosses \(c\). A discarded block remains crossing
under every coarsening. Hence this is a chain map.
The discarded summands have diagonal support outside the cut open
\(U_c\). On \(U_c\), every retained term on both sides has the
same bundle \(\mathcal K^l_{\mathbf j'\mathbf j''}\).
Consequently the cut map is an isomorphism on the full algebraic
open, before radial completion.

Repeated cut projections commute, and the empty word supplies the
counit. The resulting strong cut coalgebra has an associative
cobar for \(\star_{\mathcal K}\).
Its positive-cut contraction uses the signs and finite sums already
given after Theorem~\ref{v4-arith:asi:all-arity}.
On suspended singleton inputs the transferred operation is
\[
 B_n(sx_1,\ldots,sx_n)=s\,r_{1n}(x_1,\ldots,x_n),\qquad n\ge2;
\]
in particular the ternary operation retains the
nonconsecutive residue \(r_{13}\).

The radial frame identifies this cut map with the fusion-channel
cut after completion. Horizontal frames of two iterated cuts differ
by the integrable channel associator. Thus an identity associator
on whole-word coefficient diagrams and a nontrivial channel
associator describe the same maps in different presentations.

\subsection{Algebraic reorderings and their monodromy}
\label{v4-arith:asi:reordering}

Let \(\sigma,\tau\) be two orders of the insertion labels.
For a partition with \(r\) blocks, let \(w\) carry its root order
for \(\sigma\) to its root order for \(\tau\).
Write \(c_w\) for the ordinary permutation of the lowest-state
tensor factors. The transition is
\begin{equation}\label{v4-arith:asi:order-map}
 f\otimes e\otimes x_1\otimes\cdots\otimes x_r
 \longmapsto
 \operatorname{sgn}(w)\,
 f\otimes e\,c_w^{-1}\otimes
 x_{w^{-1}(1)}\otimes\cdots\otimes x_{w^{-1}(r)}.
\end{equation}
The sign is the determinant sign of the shifted block factors.

The Casimir connection and infinity null equation are equivariant
for these algebraic reorderings.
On the chart overlap, a partition containing two labels ordered
differently has empty support. For every surviving collision the
two colliding roots keep their relative order.
Its coordinate, integrable Ward germ and full principal part
therefore agree under \eqref{v4-arith:asi:order-map}.
If the deleted root moves from place \(i\) to place \(p\),
omitting it changes the permutation sign by \((-1)^{i+p}\).
This converts the face sign \((-1)^{i-1}\) to \((-1)^{p-1}\).

Hence the order maps commute with the differential.
They satisfy the cocycle identity by permutation coherence.
For an ordered collision, only saturated partitions survive
extraordinary pullback, and their root orders agree.
Thus the transition also commutes with every collision comparison.

On a cut-admissible partition, reordering either side is
\(w_A\times w_B\), with sign
\(\operatorname{sgn}(w_A)\operatorname{sgn}(w_B)\).
The order maps therefore glue the coupled products and their
cut isomorphisms. Interchanging the two groups gives a symmetry
whose sign on \(r_A,r_B\) blocks is \((-1)^{r_Ar_B}\).
The cut coalgebra is cocommutative for this reordering symmetry.
This symmetry does not remove the monodromy of its coefficient
bundle.

A half turn describes the order map in continued horizontal
frames. At a pair of roots, continuation of \(u^a\) through a
counterclockwise semicircle contributes \(e^{\pi ia}\).
The normalized integral-power germ changes by translation and
ordinary reordering. The exponent is retained by the coefficient
connection. Clockwise continuation gives \(e^{-\pi ia}\).
Their ratio is the full turn \(e^{2\pi ia}\).

These matrices do not generally define transitions on coefficients
at the same point for the fixed product
\eqref{v4-arith:asi:coupled-product}. The distinction already
appears for three defining \(\mathfrak{sl}_2\) colors.

\begin{proposition}\label{v4-arith:asi:half-turn-obstruction}
At an integral level \(k\ge2\), the block bundle with three
defining inputs and defining output is irreducible of rank two.
Consider a flat coefficient transition tensored with state
reordering. Compatibility with the fixed coupled product forces
one scalar on both binary channels relative to ordinary reordering.
The half turn has different scalars and cannot be that transition.
\end{proposition}

\begin{proof}
Put \(h=(k+2)^{-1}\). Both singlet and triplet channels are
integrable, so the rank is two. On the slice \((0,x,1)\),
the residues at \(0,1,\infty\) have eigenvalues
\[
 \{-3h/2,h/2\},\qquad
 \{-3h/2,h/2\},\qquad
 \{0,2h\}.
\]
All exponent differences are nonintegral.

An invariant line would extend to a line bundle whose degree is
minus a sum of residue eigenvalues. This sum is one of
\(-3h,-h,h,3h\), none an integer when \(k+2\ge4\).
Thus no invariant line exists.

A flat endomorphism of this irreducible complex local system is
scalar. Compatibility with the cut \(12|3\) extends the binary
transition over the third defining input. The two channels
therefore have the same scalar relative to reordering.
The half turn instead has eigenvalues
\[
 -e^{-3\pi ih/2},\qquad e^{\pi ih/2}.
\]
After division by the reordering signs, their ratio is
\(e^{-2\pi ih}\ne1\). This contradicts the necessary equality.
\end{proof}

At level one only the singlet channel remains.
The same rank-two argument applies at irrational \(h\) to the
full defining-color KZ system: the four possible residue sums
are again nonintegral. The proposition fixes the coefficient
transition and the coupled product. It makes no assertion
about a different tensor product with separately constructed
comparison maps.

The rank-two example also distinguishes local meridians from
full monodromy. At level two put \(\zeta=e^{\pi i/4}\).
The two local full turns can be represented as
\[
 A=\zeta Z,\qquad B=\zeta X,\qquad
 Z^2=X^2=1,\quad ZX=-XZ.
\]
Each has order eight and \(ABA^{-1}B^{-1}=-1\).
The group they generate on the slice has order sixteen:
\[
 \{\zeta^aZ^bX^c:
       a\in\mathbb Z/8,\ b,c\in\mathbb Z/2,\
       a\equiv b+c\pmod2\}.
\]

The full twist of the three moving points contributes
\(\zeta^{-3}1\), so the full pure-braid image has order thirty-two.
Projectively it is the Klein four-group.

Indeed, the exponent differences on the three-punctured sphere
are \(1/2\). Its projective local monodromies are distinct
involutions with involutive product. Irreducibility forces
their matrix lifts to anticommute.

Adjoining eighth roots of all differences kills the individual
local meridians. That abelian cover still contains their
commutator loop, whose action is \(-1\).
The cover corresponding to the kernel of the full monodromy
trivializes the bundle. On it the continued horizontal frames
are single valued and the reorderings are constant half-turn
matrices. Killing local eigenvalues alone would miss this
remaining sign.

\begin{proposition}\label{v4-arith:asi:e6-geometric-comparison}
The full ambient affine coefficient system gives a strong cut coalgebra
with the terms \eqref{v4-arith:asi:affine-cousin-term} and the coupled cut
product. It glues on the ordered collision schemes. Mark it with the
$A_E$-module junctions on its sewing charts. The arithmetic junction
comparison \eqref{v4-arith:asi:arith-affine-field} identifies these marked
residue diagrams, their cut maps, and their cobars. It retains every
normal derivative and the vacuum.
\end{proposition}

\begin{proof}
An affine Ward identity moves a negative current mode at one insertion
to regular expansions at the others. Induction on total degree
therefore determines every germ from its lowest-space value. The
singular-vector equation \eqref{v4-arith:asi:block-null} is exactly the
condition that this germ descend to the integrable quotients. At a
collision, inserting a binary integrable germ into an external germ
gives its factorization over the allowed intermediate labels. The
factorization is an isomorphism by the integrable conformal-block
sewing and factorization theorems of Tsuchiya--Ueno--Yamada,
Theorems~6.2.1 and~6.2.5.
It uses the full tensor
\eqref{v4-arith:asi:affine-propagator}.

For two disjoint collisions the two principal-part maps commute before
their determinant signs are inserted, so their signed composites
cancel. For three blocks the three routes are the three expansions of
one Ward germ. The two-variable Cousin boundary identity therefore
gives their signed sum zero. These statements prove that the sum of
the elementary faces squares to zero. They include a collision of the
first and third blocks with the second retained as a spectator.

For a consecutive cut, discard exactly the partitions having a block
on both sides. Such a block remains crossing under every coarsening,
so the projection is a chain map. Its discarded terms have diagonal
support outside the cut open. On every retained term the coupled
product has the coefficient bundle of the same whole word. The cut
map is consequently an isomorphism there. Repeated projections
commute and the empty word is the counit.

On two ordered charts, reorder the lowest-space tensor factors and
precompose their coefficient functionals by the inverse permutation.
The Casimir connection is equivariant for these algebraic reorderings.
If a colliding pair has the same relative order in both charts, its
Ward germ and principal part are carried to each other. If its order
is reversed, its diagonal is absent from the overlap. Thus the
reorderings give a cocycle of complexes.

In the continued horizontal frames of
Section~\ref{v4-arith:asi:reordering}, these algebraic maps are
represented by half-turn matrices. They are not pointwise
coefficient transitions for the fixed coupled product.

Finally \eqref{v4-arith:asi:arith-affine-field} is an equality of full
intertwining operators. Its coefficient symbols transform by the dual
inverse of $K$, preserving their contraction with fields. It therefore
commutes with \eqref{v4-arith:asi:affine-residue} and with every spectator
extension. The cut maps are the same projections on both sides.
The resulting identification of coalgebras identifies their cobars.
\end{proof}

The ambient coefficient bundle is essential here. A projection onto a
chosen fusion channel in a sewing frame need not extend to a
single-valued flat projection on the whole configuration space.
The construction retains that bundle and marks the module junctions
on the sewing charts. It does not replace the bundle by constant
channel spaces or assert a global splitting into them.

The Heisenberg, lattice, and affine constructions are distinct conformal
theories, with central charges $1$, $8$, and $5/2$. Their full state
spaces are not identified. They realize the same specified junctions
on different conformal carriers. The affine construction also supplies
nontrivial conformal-block monodromy at the embedding $\sigma$.

\section{The original embedding in a Virasoro theory}
\label{v4-arith:asi:sec:native-virasoro}

The original arithmetic embedding has $d=1-\sqrt3$. It therefore
requires a different conformal carrier from the positive affine theory.
Such a carrier is provided by the Virasoro minimal model
\[
 W_-=\operatorname{Vir}(12,5),\qquad
 c_-=1-\frac{6(12-5)^2}{12\cdot5}=-\frac{39}{10}.
\]


Let $U_j$ be its simple module with Kac label $(j+1,1)$,
$0\leq j\leq10$. Its lowest conformal weight is
\begin{equation}\label{v4-arith:asi:virasoro-weights}
 h_j^- =\frac{(5(j+1)-12)^2-49}{240}
       =\frac{5j(j+2)}{48}-\frac j2.
\end{equation}
Its states are the negative Virasoro modes on a lowest vector, modulo
the two singular-vector submodules at levels $j+1$ and $4(11-j)$.
This specifies all descendants, including their null relations.
The category is rational and \(C_2\)-cofinite.

Huang's Theorems~2.1 and~3.8--3.10 give its intertwining operators
the convergent vertex tensor structure.
Creutzig--Huang--Yang, Section~7, give the related admissible
tensor-category formulation.

Put $q_-=e^{5\pi i/12}$. The cyclotomic automorphism
$\gamma:q\mapsto q_-$ satisfies
$\gamma(\sqrt3)=-\sqrt3$ and $\gamma(i)=i$.
Write $\mathcal D_{q_-}^+$ for the even quantum-group category.
There is a ribbon equivalence
\begin{equation}\label{v4-arith:asi:virasoro-ribbon}
 \Psi:\mathcal D_{q_-}^+\xrightarrow{\ \sim\ }
                   \langle U_0,U_2,U_4,U_6,U_8,U_{10}\rangle,
 \qquad V_j\longmapsto U_j.
\end{equation}
Here and below the rational Racah bases are transported through
$\Psi$ to normalize the Virasoro intertwining operators.

We explain why the parameter and the ribbon structure in this
comparison are fixed. The fusion rules of $U_0,\ldots,U_{10}$ are
\eqref{v4-arith:asi:su2-fusion}. The object $U_1$ generates them and
$U_1^2=U_0\oplus U_2$. Its counterclockwise exchanges on these two
channels, computed from their leading primary powers, are
\[
 e^{\pi i(h_0^--2h_1^-)}=e^{3\pi i/8}=-q_-^{-3/2},\qquad
 e^{\pi i(h_2^--2h_1^-)}=e^{5\pi i/24}=q_-^{1/2}.
\]


Categories with \(SL_2\) Verlinde rules are reconstructed from
the quantum-group categories, up to a cocycle on the parity grading.
This is the classification in Etingof--Gelaki--Nikshych--Ostrik,
Section~8.27.5.

In the Temperley--Lieb reconstruction,
these two eigenvalues fix the braid generator. The successive
Jones--Wenzl projections use
\[
 [n]_- =\frac{\sin(5n\pi/12)}{\sin(5\pi/12)},
 \qquad [n]_-\ne0\ (1\leq n\leq11),\quad [12]_-=0.
\]


Their path matrix units give the $A_{11}$ fusion paths. The next
projection vanishes. Thus the comparison is full on the generated
category. A possible parity cocycle restricts to the constant one on the even
subcategory. There the twists agree exactly, since
$e^{2\pi i h_j^-}=q_-^{j(j+2)/2}$ for even $j$.
This proves \eqref{v4-arith:asi:virasoro-ribbon}, including its braiding
and twists.

In particular it fixes the orientation of analytic
continuation, which the fusion rules alone would not determine.

Apply $\gamma$ to the algebra and junction matrices of
Section~\ref{v4-arith:asi:sec:e6-tensor}. In $\mathcal D_{q_-}^+$ the
algebra is again $V_0\oplus V_6$, now with
\[
 m_{66}^{\,0}=-\frac13-\frac{\sqrt3}{6},\qquad m_{66}^{\,6}=1,
 \qquad \dim A_-=3-\sqrt3.
\]
Its associativity, commutativity, separability, and all junction
comparison equations are polynomial identities over the cyclotomic
field, so they are preserved by $\gamma$. The comparison from the
original arithmetic category has binary matrices
\begin{equation}\label{v4-arith:asi:native-junction-map}
 K^-_{ab;c}=\gamma(K_{ab;c}).
\end{equation}
Indeed $\gamma$ restricts to $\sigma$ on $F$, and applying it to
$\mathcal C^\sigma$ restores the original embedding.

Transport $A_-$ and its three modules through $\Psi$. Their full
Virasoro state spaces are
\begin{equation}\label{v4-arith:asi:native-full-states}
 \begin{split}
 \mathsf H^-_1&=U_0\oplus U_6,\\
 \mathsf H^-_g&=U_4\oplus U_{10},\\
 \mathsf H^-_\rho&=U_2\oplus U_4\oplus U_6\oplus U_8.
 \end{split}
\end{equation}
Let $q$ be the balancing quotient defined by the conjugate of
\eqref{v4-arith:asi:balancing-idempotent}.
For a module junction $p$, its full field is the intertwining operator
of the module map $\Psi(pq)\Psi_2$.
The arithmetic field is $K^-$ times this column of fields.

\begin{theorem}\label{v4-arith:asi:native-e6-comparison}
Equations \eqref{v4-arith:asi:native-junction-map} and
\eqref{v4-arith:asi:native-full-states} give an analytic ordered realization
of the order-two arithmetic category at its original complex embedding.
It has central charge $-39/10$ and relative dimensions $1,1,1-\sqrt3$.
The comparison preserves every junction, all descendants, translation,
relative duality, and collision and sewing operations in every arity.
Its sewings on the domain of
Section~\ref{v4-arith:asi:sewing-domain} converge absolutely and locally
uniformly, including finite simultaneous sewings that create handles.
\end{theorem}
\begin{proof}
The ribbon equivalence \eqref{v4-arith:asi:virasoro-ribbon} identifies the
commutative separable algebra and its balanced module tensor product.
The conjugate of \eqref{v4-arith:asi:all-f-comparison} gives all its binary
and ternary comparison maps. The proof of
Theorem~\ref{v4-arith:asi:e6-analytic-comparison} now applies with the
Virasoro vertex tensor structure in place of the affine one.
It is an equality of intertwining operators on entire modules,
so it includes their singular-vector quotients and every descendant.

The module dimensions can be checked without a choice of junction
basis. The categorical dimension of $U_j$ is
\[
 (-1)^j\frac{\sin(5(j+1)\pi/12)}{\sin(5\pi/12)}.
\]
Divide the sums for the three modules in
\eqref{v4-arith:asi:native-full-states} by $3-\sqrt3$.
They give $1,1,1-\sqrt3$. The unique spherical structure on the
near-group category is therefore the original arithmetic one.
In particular the closed marked junction has value $-e_1$.

The sewing tensor on $U_j$ is
\[
 \sum_{n\geq0}t^{h_j^-+n}
         \sum_\alpha v_{j,n,\alpha}\otimes v_{j,n,\alpha}^{\vee}.
\]
Every grade is finite dimensional. The possible negative lowest
weights contribute only fixed powers of $t$ on a chosen logarithmic
branch. Rationality and $C_2$-cofiniteness allow
Gui, Theorem 13.1 to be applied with the hypotheses of
Section~\ref{v4-arith:asi:sewing-domain}. This gives the asserted
convergence on that domain.

On a binary collision divide the field by
$u^{h_k^--h_i^--h_j^-}$ and pair it with the inverse coefficient
symbol before taking its full principal part. The derivative of that
symbol retains the conformal exponent, as in
\eqref{v4-arith:asi:affine-residue}. The full intertwiner identity commutes
with this operation and with every spectator insertion. The
Virasoro null-vector equations supply its conformal-block connection.
Their sewing identities give the same Cousin boundary and cut
identities. Thus the comparison also preserves the collision diagrams.
\end{proof}

For even $j=0,2,4,6,8,10$, the lowest weights are respectively
$0,-1/6,1/2,2,13/3,15/2$. These weights and the central charge belong
to the new Virasoro construction. The Galois
automorphism acts on the finite junction data. The independently
defined Virasoro vertex tensor theory supplies the full analytic lift.
These sewings use algebraic contragredients. No positive Hilbert
metric is assumed.

\section{The connected realization beyond order two}
\label{v4-arith:asi:connected-question}

The relative \(E_8\) construction works for every arithmetic category
considered here. Its vacuum algebra has \(N+1\) idempotents.
The affine and Virasoro constructions give complete order-two
realizations on independently defined conformal state spaces.
These results leave a different first example.

\begin{question}\label{v4-arith:asi:order-three}
For \(K=\mathbb Q(\sqrt{21})\) and
\(\varepsilon=(5+\sqrt{21})/2\), one has \(N=3\).
Construct a connected rational nonlocal conformal realization of
\(\mathcal C_I\), with its specified module fusion, marked junctions,
duality and convergent sewing on full descendant spaces.
Which rooted, collision or nearby-cycle tensor realizes its fusion?
\end{question}

Theorem~\ref{v4-arith:asi:empty-obstruction} excludes the stated
strong unital realization by the scalar-empty balanced cut tensor.
It does not exclude the alternative tensor asked for here.
Likewise, the supported recoupling comparison must retain its larger
channel system and specify the compression that produces
\(-N\beta D\). Its first missing map is an operation-preserving
supported comparison with that declared compression.

Arithmetic field automorphisms act on the finite junction
coefficients. Transport of a conformal theory requires, in addition,
its actual tensorator, state spaces and analytic sewing.
This distinction is already visible in the different affine and
Virasoro conformal weights.

\section{References}

\begin{itemize}

\item A.~Beilinson and V.~Drinfeld. \emph{Chiral Algebras}. AMS (2004), Section~3.4.11.

\item T.~Hayashi. \emph{A canonical Tannaka duality for finite semisimple tensor categories}. \href{https://arxiv.org/abs/math/9904073}{arXiv:math/9904073}.

\item C.~Dong and G.~Mason. \emph{Holomorphic vertex operator algebras of small central charges}. Pacific J. Math. \textbf{213} (2004), 253--266.

\item B.~Gui. \emph{Convergence of sewing conformal blocks}. Comm. Contemp. Math. \textbf{26} (2024), 2350007, \href{https://arxiv.org/abs/2011.07450}{arXiv:2011.07450}, Section~13 and Theorem~13.1.

\item T.~Abe, G.~Buhl and C.~Dong. \emph{Rationality, regularity, and \(C_2\)-cofiniteness}. \href{https://arxiv.org/abs/math/0204021}{arXiv:math/0204021}.

\item C.~Dong, H.~Li and G.~Mason. \emph{Regularity of rational vertex operator algebras}. \href{https://arxiv.org/abs/q-alg/9508018}{arXiv:q-alg/9508018}.

\item D.~Radchenko and C.~Wheeler. \emph{Real quadratic fields and finite quantum dilogarithms I}. \href{https://arxiv.org/abs/2609.21892v2}{arXiv:2609.21892v2}.

\item J.~Thornton. \emph{On braided near-group categories}. \href{https://arxiv.org/abs/1102.4640}{arXiv:1102.4640}.

\item A.~Kirillov, Jr., and V.~Ostrik. \emph{On a \(q\)-analogue of the McKay correspondence and the ADE classification of \(\widehat{\mathfrak{sl}}_2\) conformal field theories}. Adv. Math. \textbf{171} (2002), 183--227, \href{https://arxiv.org/abs/math/0101219}{arXiv:math/0101219}.

\item Y.-Z.~Huang and J.~Lepowsky. \emph{Intertwining operator algebras and vertex tensor categories for affine Lie algebras}. Duke Math. J. \textbf{99} (1999), 113--134, \href{https://arxiv.org/abs/q-alg/9706028}{arXiv:q-alg/9706028}.

\item M.~Finkelberg. \emph{An equivalence of fusion categories}. Geom. Funct. Anal. \textbf{6} (1996), 249--267.

\item M.~Finkelberg. \emph{Erratum to: An equivalence of fusion categories}. Geom. Funct. Anal. \textbf{23} (2013), 810--811, \href{https://doi.org/10.1007/s00039-013-0230-y}{doi:10.1007/s00039-013-0230-y}.

\item Y.-Z.~Huang and L.~Kong. \emph{Open-string vertex algebras, tensor categories and operads}. Comm. Math. Phys. \textbf{250} (2004), 433--471, \href{https://arxiv.org/abs/math/0308248}{arXiv:math/0308248}.

\item A.~Beauville. \emph{Conformal blocks, fusion rules and the Verlinde formula}. \href{https://arxiv.org/abs/alg-geom/9405001}{arXiv:alg-geom/9405001}, Sections~3--4.

\item A.~Tsuchiya, K.~Ueno and Y.~Yamada. \emph{Conformal field theory on universal family of stable curves with gauge symmetries}. Adv. Stud. Pure Math. \textbf{19} (1989), 459--566.

\item Y.-Z.~Huang. \emph{Virasoro vertex operator algebras, the (nonmeromorphic) operator product expansion and the tensor product theory}. J. Algebra \textbf{182} (1996), 201--234, \href{https://arxiv.org/abs/q-alg/9505020}{arXiv:q-alg/9505020}.

\item T.~Creutzig, Y.-Z.~Huang and J.~Yang. \emph{Braided tensor categories of admissible modules for affine Lie algebras}. Comm. Math. Phys. \textbf{362} (2018), 827--854, \href{https://arxiv.org/abs/1709.01865}{arXiv:1709.01865}.

\item P.~Etingof, S.~Gelaki, D.~Nikshych and V.~Ostrik. \emph{Tensor Categories}. AMS (2015), Section~8.27.5, p.~266.

\end{itemize}
